\section{IX. 법과 도덕 (Laws and
Morals)}\label{ix.-uxbc95uxacfc-uxb3c4uxb355-laws-and-morals}

\subsection{1. 자연법과 법실증주의 (Natural Law and Legal
Positivism)}\label{uxc790uxc5f0uxbc95uxacfc-uxbc95uxc2e4uxc99duxc8fcuxc758-natural-law-and-legal-positivism}

법과 도덕 사이에는 다양한 유형의 관계가 존재하며, 그 중 어느 하나를
\emphksans{그들 사이의} 관계로 특정하여 연구 대상으로 삼는 것은 유익하지
않다. 대신, 법과 도덕이 관련되어 있다고 주장하거나 부인할 때 의미될 수
있는 다양한 사태들을 구분하는 것이 중요하다. 때로는 거의 누구도 부인한
적 없는 일종의 연결성을 주장하는 경우도 있으며, 이러한 논박의 여지 없는
사실이 보다 의심스러운 연결성의 징표로 잘못 수용되거나, 심지어 그것과
혼동되기도 한다. 예컨대, 언제 어디서나 법의 발전이 특정 사회 집단의
관행적 도덕(moral convention)과 이상(ideal), 그리고 당시 통용되는 도덕을
초월한 개인들의 계몽된 도덕적 비판(enlightened moral criticism)에 의해
심대한 영향을 받아왔다는 점은 진지하게 반박될 수 없다. 그러나 이 진리를
부당하게 이용하여 전혀 다른 명제를 정당화하는 근거로 삼는 것도 가능하다:
즉, 법체계는 도덕 혹은 정의와의 어떤 특정한 준수/합치(a specific
conformity)를 반드시 드러\emphksans{내어야 하며}(\emph{must} exhibit), ~또는
그에 복종할 도덕적 의무가 있다는 널리 퍼진 확신(widely diffused
conviction)위에 반드시 기초\emphksans{해야 한다}(\emph{must} rest on)는
명제이다. 물론 이 주장은 어떤 의미에서는 참일 수도 있으나, 특정 법체계
내에서 사용되는 개별 법률의 법적 유효성의 기준이 반드시 명시적으로든
암묵적으로든 도덕 혹은 정의에 대한 참조(a reference to)를 포함해야
한다는 점은 이로부터 도출되지 않는다.

법과 도덕의 관계에 관해서는 이 밖에도 많은 다른 문제들이 제기될 수 있다.
이 장에서는 그 가운데 두 가지만을 논의할 것인데, 비록 두 문제 모두 다른
여러 문제들에 대한 일정한 고찰을 수반하겠지만, 여기서는 그 두 문제에
한정한다. 첫 번째는, 오늘날에는 법과 도덕에 관한 서로 다른 여러 논제를
가리키는 명칭으로 사용되게 되었음에도 불구하고, 여전히 자연법(Natural
Law)과 법실증주의(Legal Positivism) 사이의 쟁점으로 묘사하는 것이 유익할
수 있는 문제이다. 여기서 우리는 법실증주의를, \textbf{(p.~186)} 법이
일정한 도덕적 요구를 재현하거나 충족해야 한다는 것이 어떠한 의미에서도
필연적 진리(necessary truth)는 아니라는 단순한 주장으로 이해하겠다. 물론
실제로는 법이 그러한 요구를 충족해 온 경우가 종종 있었지만 말이다.
그러나 이러한 견해를 취해 온 이들이 도덕의 성격에 관해서는 침묵했거나,
또는 서로 크게 상이한 견해를 제시해 왔기 때문에, 법실증주의가 거부되어
온 두 가지 매우 상이한 형태를 검토할 필요가 있다. 그중 하나는 자연법의
고전적 이론들에서 가장 분명하게 표현되어 있는데, 이는 인간 이성이 발견을
기다리고 있는 일정한 인간 행위의 원리들이 존재하며, 인간이 제정한
법(man-made law)이 유효성(validity)을 가지기 위해서는 그러한 원리들에
합치해야 한다는 주장이다. 다른 하나는 도덕에 대해 이와는 다른, 덜
합리주의적인 관점을 취하면서, 법적 유효성이 도덕적 가치와 연결되는
방식에 대해 상이한 설명을 제시한다. 이 절과 다음 절에서는 이 가운데 첫
번째 입장을 검토할 것이다.

플라톤(Plato)에서 오늘날에 이르기까지, 인간이 어떻게 행위해야 하는지가
인간 이성(human reason)에 의해 발견될 수 있다는 명제의 언명(assertion)과
그 부정을 위해 헌신된 방대한 문헌 속에서, 논쟁의 한쪽에 선 이들은 다른
쪽에 대해 ``이것을 보지 못한다면 당신들은 눈이 멀었다''고 말하는 듯
보이며, 이에 대해 상대방은 ``당신들은 꿈을 꾸고 있다''고 응답하는 것처럼
보인다. 이러한 상황이 발생하는 이유는, 올바른 행위의 참된 원리들이
이성적으로 발견 가능하다는 주장이 보통 독립된 하나의 교리로 제시된 것이
아니라, 무생물과 생물을 포괄하는 자연(nature)에 대한 일반적 관념의
일부로서 처음 제시되었고, 오랫동안 그 일부로서 옹호되어 왔기 때문이다.
이러한 관점은 여러 측면에서 현대 세속적 사유(modern secular thought)의
틀을 구성하는 일반적 자연관과 정면으로 대립한다. 그 결과 자연법
이론(Natural Law theory)은 비판자들에게는 현대 사유가 이미 승리적으로
벗어났다고 여겨지는 깊고 오래된 혼란들로부터 비롯된 것으로 보이는 반면,
옹호자들에게는 비판자들이 피상적인 사소함에만 집착하면서 보다 심원한
진리들을 무시하고 있는 것처럼 보이게 된다.

따라서 많은 현대의 비판자들은, 적절한 행위의 법칙(laws of proper
conduct)이 인간 이성(human reason)에 의해 발견될 수 있다는 주장이
`법(law)'이라는 단어의 단순한 애매성(ambiguity)에 근거하고 있으며, 이
애매성이 드러나는 순간 자연법(Natural Law)은 치명타를 입었다고 생각해
왔다. 존 스튜어트 밀(John Stuart Mill)이 『법의 정신(Esprit des Lois)』
제1장에서, 별과 같은 무생물이나 동물들이 `그들의 자연의 법(law of their
nature)'에 복종하는데 왜 인간은 그러하지 않고 죄(sin)에 빠지는가를
순진하게 묻는 몽테스키외(Montesquieu)를 다룬 방식이 바로 그러하다. 밀은
이것이 \textbf{(p.~187)} 자연의 진행이나 정기성(regularities)을
정식화하는 법칙들과, 인간에게 특정한 방식으로 행동할 것을 요구하는
법칙들 사이의 오래된 혼동을 드러낸다고 보았다. 전자는 관찰과 추론에 의해
발견될 수 있으며 `기술적(descriptive)'이라고 불릴 수 있고, 이를 발견하는
것은 과학자의 역할이다. 반면 후자는 그러한 방식으로 확립될 수 없는데,
그것들은 사실에 대한 진술이나 기술이 아니라 인간이 특정한 방식으로
행동해야 한다는 `처방(prescriptions)' 또는 요구이기 때문이다. 따라서
몽테스키외의 질문에 대한 대답은 단순하다. 처방적 법칙(prescriptive
laws)은 위반될 수 있으며, 그럼에도 여전히 법칙으로 남아 있는데, 이는
단지 인간이 지시받은 대로 행동하지 않았다는 것을 의미할 뿐이다. 그러나
과학에 의해 발견된 자연법칙들에 대해 그것들이 위반될 수 있는지 여부를
말하는 것은 무의미하다. 만약 별들이 그들의 규칙적인 운동을 기술한다고
주장하는 과학적 법칙과는 다른 방식으로 행동한다면, 그 법칙들은 `위반된'
것이 아니라, `법칙'이라 불릴 자격을 상실하고 다시 정식화되어야 한다.
이러한 `법(law)'이라는 말의 의미 차이에 대응하여, `must', `bound to',
`ought', 'should'와 같은 관련 어휘들에서도 체계적인 차이가 나타난다.
그러므로 이러한 관점에서 보면, 자연법에 대한 신념은 매우 단순한 오류로
환원된다. 즉, 법적 함축을 지닌 이러한 단어들이 가질 수 있는 매우 상이한
의미들을 인식하지 못한 데서 비롯된 오류라는 것이다. 이는 마치 어떤
사람이 ``당신은 군복무를 위해 신고해야 한다(You are bound to report for
military service)''와 ``바람이 북쪽으로 돌면 얼게 될 것이 확실하다(It is
bound to freeze if the wind goes round to the north)''라는 문장에서 해당
표현들이 지니는 전혀 다른 의미를 인식하지 못한 것과 같은 경우이다.

벤담(Bentham)과 밀(Mill)과 같은 비판자들, 즉 자연법(Natural Law)을 가장
격렬하게 공격한 이들은, 법(law)의 이러한 상이한 의미들 사이에서
상대방들이 보인 혼동을 자연의 관찰된 정기성(regularities)이 우주의 신적
통치자(Divine Governor)에 의해 처방되거나 명령되었다고 믿는 신념의
잔존에 돌리는 경우가 많았다. 이러한 신정적(theocratic) 관점에서는,
중력의 법칙과 십계명---인간을 위한 신의 법(God's law for Man)---사이의
유일한 차이는, 블랙스톤(Blackstone)이 언명했듯이, 창조된 존재들 가운데
오직 인간만이 이성과 자유의지를 부여받았다는 비교적 사소한 점에
불과했다. 그 결과 인간은 사물들과 달리 신적 처방(divine prescriptions)을
인식할 수 있고 또한 그것을 불복종할 수도 있다는 것이다. 그러나 자연법은
항상 우주의 신적 통치자나 입법자(Divine Governor or Lawgiver)에 대한
신념과 결부되어 온 것은 아니며, 설령 그러한 신념과 결부된 경우라
하더라도 그 고유한 신조(tenets)들이 논리적으로 그 신념에 의존해 온 것은
아니다. 자연법에 포함되는 '자연적(natural)'이라는 단어의 관련된 의미와,
\textbf{(p.~188)} 처방적 법(prescriptive laws)과 기술적 법(descriptive
laws) 사이의 차이를---현대인의 눈에는 매우 명백하고
중요한---최소화하려는 그 일반적 관점은, 이러한 목적과 관련해서는
전적으로 세속적이었던 그리스 사상(Greek thought)에 그 뿌리를 두고 있다.
실제로 어떤 형태의 자연법 교설이 지속적으로 재언명되어 온 이유는
부분적으로, 그것의 호소력이 신적 권위와 인간적 권위 모두로부터
독립적이라는 점, 그리고 오늘날에는 거의 받아들여질 수 없는 용어법과 많은
형이상학을 수반하고 있음에도 불구하고, 도덕성과 법을 이해하는 데 중요한
몇 가지 기본적 진실(elementary truths)을 포함하고 있다는 점에 있다.
우리는 이러한 진실들을 그 형이상학적 틀로부터 분리해 내어, 보다 단순한
용어로 여기에서 다시 진술하고자 한다.

현대의 세속적 사유에 따르면, 무생물과 생물, 동물과 인간의 세계는 일정한
유형의 사건들과 변화들이 반복적으로 발생하며, 그 안에서 일정한 정기적
연결관계(regular connections)가 예시되는 장면이다. 이러한 연결관계들
가운데 적어도 일부는 인간에 의해 발견되고 자연법칙(laws of nature)으로
정식화되었다. 이러한 현대적 관점에서 자연을 이해한다는 것은, 자연의 어떤
부분에 대하여 이러한 정기성(regularities)에 관한 지식을 적용하는 것을
의미한다. 물론 위대한 과학 이론들의 구조가 관찰 가능한 사실, 사건, 또는
변화를 어떤 단순한 방식으로 그대로 반영하는 것은 아니다. 실제로 그러한
이론들의 상당 부분은 관찰 가능한 사실과 직접적으로 대응되는 대상이 없는
추상적인 수학적 정식화(abstract mathematical formulations)로 이루어져
있다. 그럼에도 불구하고 이러한 이론들이 관찰 가능한 사건들과 변화들과
연결되는 지점은, 이러한 추상적 정식화들로부터 관찰 가능한 사건들을
참조하며, 그러한 사건들에 의해 확증되거나 반증될 수 있는 일반화들이
연역될 수 있다는 데에 있다. 따라서 과학 이론이 자연에 대한 우리의 이해를
증진시킨다고 언명할 수 있는 근거는, 궁극적으로는 무엇이 발생할 것인지를
예측할 수 있는 그 이론의 능력에 달려 있으며, 이 능력은 정기적으로
발생하는 것들에 대한 일반화에 기초한다. 중력의 법칙과 열역학 제2법칙은,
현대적 사유에서 볼 때, 관찰 가능한 현상들의 정기성에 관하여 제공하는
정보 덕분에, 단순한 수학적 구성물에 그치지 않고 자연법칙(laws of
nature)으로 간주된다.

자연법 교설은, 관찰 가능한 세계를 단지 그러한 정기성(regularities)의
장면으로만 보지 않고, 자연에 대한 앎을 그 정기성에 관한 앎으로만
이해하지 않는, 보다 오래된 자연 개념의 일부이다. 오히려 이러한 오래된
관점에 따르면, 인간적 존재이든, 생물적 존재이든, 무생물적 존재이든, 이름
붙일 수 있는 모든 종류의 존재하는 것은 단지 자기 자신의 존재를
유지하려는 경향을 지닐 뿐만 아니라, \textbf{(p.~189)} 그것에 고유한
선(good)---곧 그것에 적합한 \emphksans{목적(end)}---에 해당하는 일정한 최적
상태(definite optimum state)를 향해 나아가는 것으로 파악된다.

이는 자연이 그 자체 안에 사물들이 실현해 나가는 탁월성의 단계(levels of
excellence)를 포함하고 있다는 목적론적 자연관(teleological conception of
nature)이다. 어떤 특정한 종류에 속하는 사물이 그 고유하거나 적합한
목적(specific or proper end)에 이르기까지 거치는 단계들은
정기적(regular)이며, 사물의 특징적인 변화 방식, 작용 방식, 또는 발전
방식을 서술하는 일반화로 정식화될 수 있다; 이러한 점에서 목적론적
자연관은 현대적 사고와 중첩된다. 차이는, 목적론적 관점에서는 사물들에게
정기적으로 일어나는 사건들이 \emphksans{단지(merely)} 정기적으로 발생하는
것으로만 이해되지 않으며, 그것들이 \emphksans{실제로(do)} 정기적으로
발생하는지 여부, 그것들이 발생\emphksans{해야 하는지(should)} 여부, 또는
그것들이 발생하는 것이 \emphksans{좋은지(good)} 여부라는 질문들이 서로 분리된
질문으로 간주되지 않는다는 데 있다. 오히려('우연(chance)'에 귀속되는
몇몇 드문 기형적 사례들을 제외하면), 일반적으로 발생하는 것은 해당
사물의 적합한 목적이나 목표를 향해 나아가는 하나의 단계로 제시됨으로써,
설명될 수 있을 뿐만 아니라 선한 것 또는 발생해야 할 것으로 평가될 수
있다. 그러므로 사물의 발전에 관한 법칙들은, 그것이 어떻게 마땅히
행동하거나 변화해야 하는지와, 그것이 실제로 어떻게 정기적으로 행동하거나
변화하는지를 모두 보여주어야 한다.

이와 같은 자연에 대한 사고방식은 추상적으로 진술될 경우 낯설게 보인다.
그러나 오늘날에도 적어도 생명체에 관해서 우리가 여전히 사용하는 몇몇
표현 방식을 상기한다면, 그것은 덜 기이하게 보일 수도 있다. 왜냐하면
목적론적 관점(teleological view)은 여전히 그들의 발전을 서술하는
일상적인 방식들 속에 반영되어 있기 때문이다. 예컨대 도토리(acorn)의
경우, 참나무(oak)로 성장하는 것은 도토리들이 정기적으로 달성하는 어떤
상태일 뿐만 아니라, 또한 (마찬가지로 정기적으로 일어나는)
부패(decay)와는 달리, 성숙의 최적 상태(optimum state of maturity)로서
구별된다. 그리고 이 최적 상태에 비추어 중간 단계들은 설명되며, 선한지
나쁜지로 평가되고, 그 다양한 부분들의 `기능(functions)'과 구조적
변화들이 식별된다. 잎의 정상적인 성장은 `완전한(full)' 또는
`적합한(proper)' 발전에 필요한 수분을 얻기 위해 요구되며, 그러한 수분을
공급하는 것이 잎의 '기능'이다. 따라서 우리는 이러한 성장을 '자연적으로
발생해야 할 것(ought naturally to occur)'으로 사고하고 말한다. 반면
무생물의 작용이나 운동의 경우에는, 그것들이 인간에 의해 어떤 목적을 위해
설계된 인공물(artefacts)이 아닌 한, 이러한 방식의 말하기는 훨씬 설득력이
떨어진다. 돌이 땅으로 떨어질 때 어떤 적합한 \textbf{(p.~190)}
'목적(end)'을 실현하고 있거나, 마치 말이 마구간으로 달려 돌아가듯이
자신의 '고유한 자리(proper place)'로 되돌아가고 있다는 관념은,
오늘날에는 다소 희극적으로 느껴진다.

실로, 자연에 대한 목적론적 관점(teleological view of nature)을 이해하는
데에는 하나의 난점이 있는데, 그것은 이 관점이 정기적으로 일어나는 바에
대한 진술과 일어나야 할 바에 대한 진술 사이의 차이를 최소화하듯이, 또한
현대 사상에서 그토록 중요한 차이로 여겨지는, 스스로의 목적을
\emphksans{가지고(with)} 그것을 의식적으로 실현하려고 노력하는 인간과, 다른
생명체나 무생물 사이의 차이 역시 최소화하기 때문이다. 왜냐하면 세계에
대한 목적론적 관점에서는 인간 역시 다른 사물들과 마찬가지로, 그에게
설정된 특정한 최적 상태(optimum state) 또는 목적(end)을 향해 나아가는
존재로 이해되며, 그가 다른 사물들과 달리 이를 의식적으로 수행할 수
있다는 사실은 그를 자연의 나머지 부분과 근본적으로 구별하는 차이로
간주되지 않기 때문이다. 이러한 특정한 인간의 목적 또는 선(good)은
부분적으로는 다른 생명체들과 마찬가지로 생물학적 성숙의 상태와 발달된
신체적 능력의 조건을 포함하지만, 또한 사유(thought)와 행위(conduct)에
나타나는 정신과 인격의 발달과 탁월성이라는, 인간에게 특유한 요소를
포함한다. 다른 사물들과 달리 인간은 추론(reasoning)과 성찰(reflection)을
통해 이러한 정신과 인격의 탁월성을 달성하는 것이 무엇을 의미하는지를
발견하고, 그것을 욕구(desire)할 수 있다. 그럼에도 불구하고, 이 목적론적
관점에 따르면 이러한 최적 상태는 인간이 그것을 욕구하기 때문에 그의
선이나 목적이 되는 것이 아니라, 오히려 그것이 이미 그의 자연적
목적(natural end)이기 때문에 인간이 그것을 욕구하는 것이다.

다시 말해, 이러한 목적론적 관점의 상당 부분은 우리가 인간을 사고하고
말하는 몇몇 방식 속에 여전히 살아 있다. 그것은 우리가 어떤 것들을 인간의
\emphksans{욕구(needs)}로 식별하여 그것을 충족시키는 것이 \emphksans{좋은(good)}
것이라고 여기거나, 인간에게 행해지거나 인간에 \emphksans{의해(by)} 겪어지는
어떤 것들을 \emphksans{해악(harm)} 또는 \emphksans{상해(injury)}로 규정하는 데
잠재되어 있다. 따라서 일부 사람들은 죽기를 원하기 때문에 먹거나 쉬기를
거부할 수도 있다는 점이 사실이기는 하지만, 우리는 먹고 쉬는 것을 단지
인간이 정기적으로 하는 행위이거나 우연히 욕구하는 것 이상으로 생각한다.
음식과 휴식은, 어떤 이들이 그것이 필요한 때에도 이를 거부한다 하더라도,
인간의 욕구이다. 그러므로 우리는 모든 인간이 먹고 잠자는 것이
자연스럽다고 말할 뿐만 아니라, 모든 인간은 때때로 먹고 쉬어야 한다고,
혹은 이러한 행위를 하는 것이 자연적으로 좋다고 말한다. 이러한 인간
행위에 대한 판단에서 `자연적으로(naturally)'라는 말의 힘은, 그것들을
단순한 관행(convention)이나 인간의 규정(prescription)을 반영하는
판단(`모자를 벗어야 한다')---그 내용은 \textbf{(p.~191)} 사유(thought)나
성찰(reflection)에 의해 발견될 수 없는 것---과도 구별하고, 또한 어떤
특정한 목적을 달성하기 위해 요구되는 것이 무엇인지를 지시할 뿐인 판단,
즉 어느 한 시점에는 한 사람이 우연히 가질 수도 있고 다른 사람은 갖지
않을 수도 있는 목적에 관한 판단과도 구별하는 데 있다. 동일한 관점은 신체
기관의 \emphksans{기능(functions)}에 대한 우리의 개념과, 그것들을 단순한
인과적 속성(causal properties)과 구별하여 긋는 경계에도 나타난다. 우리는
심장의 기능이 혈액을 순환시키는 것이라고 말하지만, 암성 종양(cancerous
growth)의 기능이 죽음을 초래하는 것이라고 말하지는 않는다.

인간 행위에 대한 일상적 사고 속에 여전히 살아 있는 목적론적 요소들을
설명하기 위해 제시된 이러한 소박한 예들은, 인간이 다른 동물들과 공유하는
생물학적 사실이라는 낮은 영역에서 취해진 것이다. 이 사고방식과 표현을
이해 가능하게 만드는 것이 전적으로 자명한 어떤 것이라는 점은 정당하게
지적될 것이다. 그것은 곧 인간 활동의 적절한 목적(end)이
생존(survival)이라는 묵시적 가정이며, 이는 대부분의 인간이 대부분의 시간
동안 자신의 존재를 계속 유지하기를 바란다는 단순한 우연적
사실(contingent fact)에 근거한다. 우리가 자연적으로 행하는 것이 좋다고
말하는 행위들은 생존에 요구되는 것들이며, 인간의 욕구(human need),
해악(harm), 그리고 신체 기관이나 변화의 \emphksans{기능(function)}이라는
개념들도 동일한 단순한 사실에 기초한다. 물론 여기에서 멈춘다면 우리는
자연법(Natural Law)의 매우 희석된 버전만을 얻게 될 것이다. 왜냐하면 이
관점의 고전적 옹호자들은 생존(\emph{自己存在保存(perseverare in esse
suo)})을 인간의 목적(end) 또는 인간에게 선(good)으로서의 훨씬 더
복잡하고 훨씬 더 논쟁적인 개념을 이루는 가장 낮은 층위로만 이해했기
때문이다. 아리스토텔레스(Aristotle)는 여기에 인간 지성의 무욕적
계발(disinterested cultivation of the human intellect)을 포함시켰고,
아퀴나스(Aquinas)는 신에 대한 앎(knowledge of God)을 포함시켰는데, 이
둘은 모두 도전될 수 있었고 실제로 도전되어 온 가치들이다. 그러나 다른
사상가들, 그중에서도 홉스(Hobbes)와 흄(Hume)은 시야를 더 낮추는 데
기꺼이 동의해 왔다. 그들은 생존이라는 소박한 목적 속에서 자연법의
용어법에 경험적 건전성(empirical good sense)을 부여하는 중심적이고
논쟁의 여지가 없는 요소를 보았다. ``인간의 본성은 어떤 경우에도 개인들의
결합(association) 없이 존속할 수 없으며, 그러한 결합은 형평(equity)과
정의(justice)의 법칙들에 대한 고려가 전혀 없다면 결코 성립할 수 없었을
것이다.''\footnote{Hume, \emph{Treatise of Human Nature}, m. ii, `Of
  Justice and Injustice'.}

이 단순한 사유는 사실 법과 도덕 둘 모두의 성격과 매우 깊은 관련을
가지며, \textbf{(p.~192)} 인간에게 있어 목적(end) 또는 선(good)이 하나의
특정한 삶의 방식으로 제시되는 일반적인 목적론적 관점 가운데서, 실제로
사람들 사이에 깊은 불일치가 존재할 수 있는 보다 논쟁적인 부분들로부터
분리해낼 수 있다. 더 나아가 우리는 생존(survival)을 언급함에 있어,
그것이 인간의 고유한 목적이나 목표이기 때문에 인간이 필연적으로 욕구하는
것으로 선험적으로 고정되어 있다는 관념을, 현대적 사고에 비추어 지나치게
형이상학적인 것으로서 배제할 수 있다. 대신 우리는, 일반적으로 인간이
살아 있기를 욕구한다는 것이 달리 그러하지 않을 수도 있었던 하나의 우연적
사실(contingent fact)에 불과하며, 생존을 인간의 목적 또는 목표라고 부를
때에도 인간이 실제로 그것을 욕구한다는 것 이상을 의미하지 않을 수도
있다고 볼 수 있다. 그럼에도 불구하고, 우리가 이러한 상식적인 방식으로
생존을 이해하더라도, 생존은 여전히 인간의 행위와 그에 대한 우리의
사유에서 특별한 지위를 차지하며, 이는 정통적 자연법(Natural Law)의
정식화들에서 생존에 부여되는 두드러진 중요성과 필연성에 상응한다. 이는
단지 압도적 다수의 인간이 극심한 고통을 대가로 치르더라도 살아 있기를
원한다는 사실 때문만이 아니라, 이러한 사실이 우리가 세계와 서로를
기술하는 데 사용하는 사유와 언어의 전체 구조 속에 반영되어 있기
때문이다. 우리는 살아 있으려는 일반적인 욕구를 제거한 채로
위험(danger)과 안전(safety), 해악(harm)과 이익(benefit), 필요(need)와
기능(function), 질병(disease)과 치료(cure)와 같은 개념들을 온전히 유지할
수 없다. 왜냐하면 이러한 개념들은, 하나의 목표로 받아들여진 생존에
기여하는 바에 참조하여 사물을 동시에 기술하고 평가하는 방식들이기
때문이다.

그러나 이러한 것들보다 더 단순하고 덜 철학적인 고려들도 존재하는데, 이는
인간의 법과 도덕에 대한 논의와 보다 직접적으로 관련되는 의미에서, 생존을
하나의 목적(aim)으로 받아들이는 것이 필수적임을 보여준다. 우리는 논의의
용어들 자체가 전제하는 어떤 것으로서 이미 이에 관여(commit)되어 있다.
왜냐하면 우리의 관심사는 자살 클럽의 사회적 장치가 아니라, 지속적인
존재를 위한 사회적 장치들이기 때문이다. 우리는 이러한 사회적 장치들
가운데, 이성에 의해 발견 가능한 자연법(natural laws)으로서 조명될 수
있는 것들이 있는지, 그리고 그것들이 인간의 법과 도덕과 어떤 관계에
있는지를 알고자 한다. 인간들이 \emphksans{어떻게} 함께 살아야 하는지에 관한
이 문제 또는 그 밖의 어떤 문제를 제기하기 위해서는, 일반적으로 말해
그들의 목적이 살아가는 데 있다는 점을 가정하지 않으면 안 된다. 이
지점에서부터 논증은 단순하다. 인간의 본성과 인간이 살아가는 세계에 관한
몇 가지 매우 명백한 일반화들---사실상 자명한 진술들(truisms)---에 대한
성찰은, \textbf{(p.~193)} 그것들이 유효한 한, 어떤 사회 조직이 존속
가능하기 위해서는 반드시 포함해야 할 일정한 행위 규칙들이 존재함을
보여준다. 이러한 규칙들은 실제로, 사회적 통제가 서로 다른 형태로 구별될
정도로 발전한 모든 사회들의 법과 관행적 도덕(conventional morality) 속에
공통 요소를 이루고 있다. 이와 함께, 법과 도덕 양자 모두에서, 특정 사회에
고유한 많은 요소들과, 임의적이거나 단순한 선택의 문제로 보일 수 있는
많은 요소들도 발견된다. 인간, 그들의 자연적 환경, 그리고 목적에 관한
기초적 진실들에 근거를 둔 이러한 보편적으로 승인된 행위 원칙들은, 그
이름 아래 흔히 제시되어 온 보다 웅대하고 더 논쟁적인 구성들과 대비하여,
자연법(Natural Law)의 \emphksans{최소 내용(minimum content)}으로 간주될 수
있다. 다음 절에서는, 이 겸손하지만 중요한 최소 내용이 의존하고 있는 인간
본성의 두드러진 특성들을 다섯 가지 자명한 명제(truisms)의 형태로 검토할
것이다.

\subsection{2. 자연법의 최소 내용 (The Minimum Content of Natural
Law)}\label{uxc790uxc5f0uxbc95uxc758-uxcd5cuxc18c-uxb0b4uxc6a9-the-minimum-content-of-natural-law}

여기에서 제시하는 단순한 자명한 이치들(truisms)과 그것들이 법과 도덕과
맺는 연관을 고찰함에 있어, 각 경우마다 언급된 사실들이 생존을 하나의
목적(aim)으로 전제할 때 법과 도덕이 특정한 내용을 포함해야 하는
\emphksans{이유(reason)}를 제공한다는 점을 관찰하는 것이 중요하다. 논증의
일반적 형태는 단순히, 그러한 내용이 없다면 법과 도덕은 인간들이 서로
결합하여 살아가는 데서 갖는 최소한의 목적, 즉 생존이라는 목적을 증진시킬
수 없다는 것이다. 이러한 내용이 결여될 경우, 인간은 있는 그대로의
상태에서 어떠한 규칙도 자발적으로 준수할 이유를 갖지 못할 것이며, 또한
규칙에 복종하고 이를 유지하는 것이 자신의 이익에 부합한다고 판단하는
사람들에 의해 자발적으로 제공되는 최소한의 협력이 없다면, 자발적으로
따르지 않으려는 다른 사람들에 대한 강제(coercion)는 불가능할 것이다.
이러한 접근에서 자연적 사실들과 법적·도덕적 규칙의 내용 사이에 존재하는
독특하게 합리적인 연결(rational connection)을 강조하는 것은 중요하다.
왜냐하면 자연적 사실들과 법적 또는 도덕적 규칙 사이에는 전혀 다른 형태의
연결이 존재하는지에 대해서도 탐구하는 것이 가능하고 또한 중요하기
때문이다. 예컨대 아직 젊은 과학인 심리학과 사회학은, 일정한
물리적·심리적·경제적 조건들이 충족되지 않는 한---예컨대 어린 아이들이
\textbf{(p.~194)} 가족 내에서 특정한 방식으로 먹여지고 양육되지 않는
한---어떠한 법 체계나 도덕 규범도 성립할 수 없거나, 혹은 오직 특정한
유형에 부합하는 법들만이 성공적으로 기능할 수 있다는 사실을 발견할 수도
있고, 이미 발견했을 수도 있다. 이러한 종류의 자연적 조건들과 규칙 체계
사이의 연결은 \emphksans{이유(reasons)}에 의해 매개되지 않는다. 왜냐하면
그것들은 특정 규칙들의 존재를, 그 규칙들의 주체들이 의식적으로 갖는
목적이나 목표와 연관시키지 않기 때문이다. 유아기에 특정한 방식으로
양육되는 것이 어떤 인구 집단이 도덕적 또는 법적 규범을 발전시키거나
유지하는 데 필요한 조건, 심지어 하나의 \emphksans{원인(cause)}임이 입증될
수는 있겠지만, 그것이 그들이 그렇게 하는 \emphksans{이유}는 아니다. 이러한
인과적 연결들은 목적이나 의식적 목표에 기초한 연결들과 물론 충돌하지
않는다. 오히려 그것들은 자연법(Natural Law)이 출발점으로 삼는 인간의
의식적 목적이나 목표가 왜 형성되는지를 실제로 설명해 줄 수 있다는
점에서, 후자보다 더 중요하거나 더 근본적인 것으로 간주될 수도 있다.
이러한 유형의 인과적 설명들은 자명한 이치들(truisms)에 기초하지도
않으며, 의식적 목표나 목적에 의해 매개되지도 않는다. 그것들은 관찰에
기초한 일반화와 이론의 방법, 그리고 가능한 경우에는 실험에 의존하는
방식으로, 사회학이나 심리학과 같은 다른 과학들이 확립해야 할 문제이다.
그러므로 이러한 연결들은, 다음에 제시되는 자명한 명제들에 언급된
사실들과 특정한 법적·도덕적 규칙들의 내용을 연관시키는 연결과는 다른
종류의 것이다.

 (i) \emphksans{인간의 취약성} (Human vulnerability) 법과 도덕의 공통된
요구사항들은 대부분 적극적 봉사가 아니라 자제(forbearance)의 형태로
나타나며, 이는 대개 부정적 명령, 즉 금지(prohibition)로 표현된다.
사회생활에 있어 가장 중요한 것은 타인에게 신체적 위해를 가하거나
살해하는 데에 폭력을 사용하는 것을 제한하는 규칙들이다. 이러한
규칙들의 기본적 성격은 다음과 같은 질문을 통해 드러난다. 만약 이런
규칙이 없다면, 우리 같은 존재들이 \emphksans{다른 종류의} 어떤 규칙이라도
가질 이유가 있을까? 이 수사적 질문의 힘은 인간이 때때로 타인을
신체적으로 공격하는 경향이 있으며, 일반적으로는 타인의 공격에
취약하다는 사실에 기반한다. 이는 상식적 진술(truism)이지만, 필연적
진리(necessary truth)는 아니다. 상황은 달랐을 수도 있고, 언젠가는
달라질 수도 있다. 어떤 동물 종들은 외골격(exoskeleton)이나 딱딱한
등껍질(carapace)과 같은 신체 구조를 지니고 있어, 동종의 다른 개체의
공격에 사실상 면역이거나, 공격 수단이 전혀 없는 경우도 있다. 만약
인간이 \textbf{(p.~195)} 서로에게 신체적으로 취약하지 않게 된다면,
법과 도덕의 가장 대표적인 규정 가운데 하나인 \emphksans{살인하지 말라}(Thou
shalt not kill)는 조항의 명백한 이유도 사라질 것이다.

 (ii) \emphksans{근사적 동등성} (Approximate equality) 인간은 신체적 힘, 민첩성,
그리고 더욱이 지적 능력(intellectual capacity)에서 서로 다르다.
그럼에도 불구하고, 인간 법제와 도덕의 다양한 형태를 이해함에 있어
결정적으로 중요한 사실은, 어떤 개인도 타인들을 장기간에 걸쳐 협력 없이
지배하거나 제압할 수 있을 정도로 압도적으로 강하지 않다는 점이다. 가장
강한 사람조차도 때때로 잠을 자야 하며, 수면 중에는 일시적으로 그의
우월성을 상실한다. 이와 같은 \emphksans{근사적 동등성}(approximate
equality)의 사실은, 상호 자제(mutual forbearance)와 절충(compromise)의
체계가 필요하다는 점---법적·도덕적 의무의 기반이 되는 점---을 무엇보다
명확히 보여준다. 이러한 자제를 요구하는 규칙들과 함께 살아가는
사회생활은 때로는 불편하고 성가시기도 하지만, 이와 같은 대략적으로
평등한 존재들에게 있어 그것은 제한 없는 공격 상태에서의 삶보다는 덜
악랄하고, 덜 야만적이며, 덜 짧다. 물론 이와 양립하는 또 하나의 상식적
진술은, 이러한 자제 체계가 수립되면 언제나 이를 악용하려는 자들이
존재한다는 것이다. 즉, 그들은 동시에 이 체계의 보호 아래 살면서도 그
제약을 깨뜨리려 한다. 실제로 이는, 뒤에서 설명하겠지만, 단순한 도덕적
통제로부터 조직화된 법적 통제로 이행하는 것이 필연적인 이유가 되는
자연적 사실 가운데 하나이다. 다시 말해, 사태는 다르게 전개될 수도
있었을 것이다. 인간이 대략적으로 평등한 대신, 어떤 사람들은 현존 평균
수준보다 훨씬 강하거나, 휴식을 거의 취하지 않아도 되는 존재일 수도
있었고, 또는 대부분의 사람들이 평균보다 훨씬 열등했을 수도 있다.
이러한 예외적 존재들에게는, 타인과 상호 자제하거나 타협하는 것보다
공격을 통해 더 많은 것을 얻을 가능성이 컸을 것이다. 그러나 우리는
거인들이 소인들 사이에서 사는 환상적인 장면을 상상할 필요도 없다.
왜냐하면 \emphksans{근사적 동등성}의 결정적 중요성은 국제 사회의
사실들---국가 간의 힘과 취약성의 극단적 불균형---에서 오히려 더 잘
드러나기 때문이다. 후술하겠지만, 국제법(international law)의 구성
단위들 사이의 이러한 불평등은, 그것이 국내법(municipal law)과 매우
다른 성격을 갖게 만든 요인 중 하나이며, 국제법이 조직화된 강제
체계(coercive system)로 기능할 수 있는 범위를 제한해왔다.
\textbf{(p.~196)}

 (iii) \emphksans{제한된 이타성} (Limited altruism) 인간은 서로를 말살하려는
욕망에 지배당한 악마가 아니며, 단지 생존이라는 소박한 목적만을
가정하더라도, 법과 도덕의 기본 규칙들이 필수적이라는 점을 입증하는
것은, 인간이 본질적으로 이기적이며 타인의 생존과 복지에 대해
무관심하다는 잘못된 견해와 동일시되어서는 안 된다. 그러나 인간이
악마가 아니라면, 동시에 천사도 아니다. 인간이 이 두 극단 사이의
중간자라는 사실은 상호 자제 체계가 \emphksans{필요하며 가능}하다는 점을
보여준다. 타인에게 해를 끼칠 유혹조차 느끼지 않는 천사들 사이에서는,
자제를 요구하는 규칙이 필요하지 않을 것이다. 반면에, 자기 희생을
개의치 않고 파괴를 일삼는 악마들 사이에서는 그러한 규칙이 아예 성립할
수 없다. 실제 세계에서 인간의 이타성은 범위가 제한적이고 지속적이지
않으며, 공격적 성향은 충분히 자주 발현되어 통제되지 않으면 사회생활을
파괴할 수 있다.

 (iv) \emphksans{제한된 자원} (Limited resources) 인간은 음식, 의복, 거처를
필요로 하며, 이들은 무한히 풍부하게 존재하지 않고, 자연으로부터
경작하거나 획득하거나 인간의 노동으로 만들어져야 한다는 사실은 단지
우연적인(contingent) 사실일 뿐이다. 그러나 이 사실들만으로도, 어떤
최소 형태의 재산 제도(property institution)는 필수적이며(비록 그것이
반드시 개인 재산일 필요는 없더라도), 이를 존중할 것을 요구하는 고유한
규칙 유형이 필요하다. 가장 단순한 재산 개념은 다음과 같은 규칙에서
찾아볼 수 있다. 즉, `소유자'(owner)가 아닌 자들의 토지 출입 또는 사용,
혹은 물질적 사물의 취득 및 사용을 배제하는 규칙이다. 작물이 자라기
위해서는 토지가 무차별적 출입으로부터 보호되어야 하고, 식량은 수확
또는 포획과 소비 사이의 기간 동안 타인에 의해 빼앗기지 않아야 한다.
언제 어디서든, 인간의 삶은 이러한 최소한의 자제 행위에 의존한다. 물론
이 점 또한 상황이 달랐을 수도 있다. 인간 유기체가 식물처럼 대기 중에서
음식물을 추출할 수 있었거나, 인간이 필요로 하는 것이 경작 없이 무한히
자라나는 방식으로 구성되어 있었다면 말이다.

지금까지 논의한 규칙들은 \emphksans{정태적} 규칙(\emph{static} rules)이다.
즉, 그것들이 부과하는 의무 및 그 의무의 귀속은 개인에 의해 변형될 수
없는 의미에서 그러하다. 그러나 최소 규모 집단을 제외한 모든 집단은,
충분한 자원을 확보하기 위해 노동의 분업(division of labour)을 발전시켜야
하며, 이는 개별 구성원이 의무를 생성하거나 그 귀속을 변경할 수 있도록
하는 \textbf{(p.~197)} \emphksans{동태적} 규칙(\emph{dynamic} rules)을 필요로
한다. 그 예로는 생산물을 이전, 교환, 또는 판매할 수 있게 하는 규칙들이
있다. 이러한 거래들은, 재산의 가장 단순한 형태를 정의하는 초기 권리와
의무의 귀속을 변경할 수 있는 능력을 전제하기 때문이다. 동일하게 불가피한
노동 분업과 지속적인 협력의 필요성은, 약속이 의무의 원천(source of
obligation)으로 인정되도록 하는 다른 유형의 동태적 규칙들을 필수적인
것으로 만든다. 이러한 장치를 통해, 개인은 말(구두든 문서든)을 통해
자신을 일정한 방식으로 행동할 의무를 지닌 존재로 만들 수 있으며, 그 의무
불이행 시 비난이나 처벌의 대상이 될 수 있다. 이타성이 무제한이 아니라면,
그러한 자기구속적 행위(self-binding act)를 제도적으로 보장하는 절차가
요구된다. 이는 타인의 미래 행위에 대한 최소한의 신뢰를 가능케 하고,
협력에 필수적인 예측 가능성(predictability)을 확보하기 위한 것이다. 이는
특히 교환되거나 공동으로 기획되는 것이 상호 서비스이거나, 재화(goods)가
동시에 또는 즉각적으로 제공되지 않는 경우에 가장 필요하다.

 (v) \emphksans{제한된 이해력과 의지력} (Limited understanding and strength of
will) 인격(persons), 재산(property), 약속(promises)을 존중하는
규칙들이 사회생활에서 필요하다는 사실은 단순하며, 그 상호적 이익도
명백하다. 대부분의 사람들은 이를 인식할 수 있으며, 이러한 규칙을
준수하는 데 요구되는 단기적 이익의 희생을 감수할 수 있다. 실제로
사람들은 다양한 동기에 따라 복종할 수 있다. 어떤 이는 그러한 희생이
그에 상응하는 이익을 가져온다는 타산적 계산(prudential
calculation)에서, 또 어떤 이는 타인의 복지에 대한 사욕없는
관심(disinterested interest)에서, 또 어떤 이는 해당 규칙 자체가 존중할
가치가 있다고 보고, 그것에 헌신하는 데서 이상(ideals)을 찾기 때문이다.
그러한 이해에 기초한 복종 동기의 실효성(efficacy)을 좌우하는 의지의
강도나 선함은 모든 사람에게 동등하게 분포되어 있지 않다. 누구나 때때로
자신의 즉각적 이익을 우선하려는 유혹을 받으며, 이를 탐지하고 처벌하기
위한 특수한 조직이 없다면, 많은 사람들이 그 유혹에 굴복하게 될 것이다.
물론 상호 자제의 이점이 워낙 명백하기 때문에, \textbf{(p.~198)} 강제적
체계(coercive system) 내에서 자발적으로 협력하려는 사람들의 수와 그
힘은, 일반적으로 범법자들의 가능한 조합보다 클 것이다. 그럼에도
불구하고, 아주 소규모의 밀접하게 조직된 사회를 제외하고는, \emphksans{억제
체계(system of restraints)}에 대한 복종은, 그 의무(obligations)에
복종하지 않으면서 그 이점만을 얻고자 하는 자들을 강제할 조직이 없다면,
어리석은 일이 될 것이다. 따라서 `제재'(sanctions)는 통상적인 복종의
동기로 요구되는 것이 아니라, 자발적으로 복종하려는 자들이 복종하지
않으려는 자들에게 희생되지 않도록 \emphksans{보장하는 장치(guarantee)}로서
요구된다. 이러한 장치 없이 복종한다는 것은 도태될 위험을 감수하는
일이다. 이와 같은 항구적인 위험이 존재하는 조건에서, 이성이 요구하는
것은 \emphksans{강제적} 체계 내에서의 \emphksans{자발적} 협력(\emph{voluntary}
co-operation in a \emph{coercive} system)이다.

이 지점에서 주목해야 할 것은, 인간 사이의 근사적 동등성(approximate
equality)이라는 동일한 자연적 사실이, 조직화된 제재가 효과를 발휘하는 데
결정적으로 중요하다는 점이다. 만약 어떤 이들이 다른 이들보다 훨씬 더
강력하여 타인의 자제에 의존하지 않아도 된다면, 범법자들의 힘이 법과
질서를 지지하는 이들의 힘을 능가할 수 있다. 이러한 불평등이 존재할 경우,
제재의 사용은 성공할 수 없으며, 오히려 그것이 억제하려는 위험만큼이나 큰
위험을 초래할 수 있다. 이러한 조건에서는, 사회생활이 소수 범법자에 대해
간헐적으로 힘을 사용하는 상호 자제 체계에 기반하기보다는, 약자들이
가능한 최상의 조건 하에서 강자에게 굴복하고, 그들의 `보호' 하에 살아가는
체계만이 유일하게 존속 가능하다. 이러한 체계는 자원의 희소성(scarcity of
resources)으로 인해 다수의 상호 충돌하는 권력 중심들이 생기게 하며, 각
중심은 자신만의 `강자'(strong man)를 중심으로 결집하게 된다. 이들은
때때로 상호 간 전쟁을 벌일 수도 있지만, 패배라는 자연적 제재(natural
sanction)의 위험이 결코 무시할 수 없는 것이기에 불안정한 평화를 유지할
수도 있다. `권력자들'(powers)이 직접 전투를 벌이기를 꺼리는 쟁점들에
대해서는, 일종의 규칙들이 수용될 수도 있다. 다시 말해, 우리는 소인국과
거인의 환상적 비유를 상상할 필요 없이, 근사적 동등성의 단순한 작동
논리(logistics)와 그것이 법에 있어 갖는 중요성을 이해할 수 있다. 국가들
간의 힘의 격차가 극심했던 국제 사회(international scene)는 그 예시로
충분하다. 수 세기 동안 국가 간의 이러한 불균형은 조직화된 제재의 작동을
불가능하게 만들었고, 법은 `중대 사안'(vital issues)에 영향을 미치지 않는
영역에만 한정되었다. \textbf{(p.~199)} 향후 원자무기(atomic weapons)가
모든 국가에 의해 사용 가능해질 경우, 이러한 불균형이 어느 정도 시정되고,
국제 통제가 국내 형사법(municipal criminal law)에 더 가까운 형태로
진화하게 될지 여부는 아직 확정되지 않았다.

우리가 지금까지 논의한 이러한 단순한 자명한 이치들(truisms)은,
자연법(Natural Law) 이론 속의 핵심적인 이성적 직관을 드러낼 뿐 아니라,
법과 도덕을 이해하는 데 본질적으로 중요하며, 왜 법과 도덕의 기본 형태를
특정한 내용이나 사회적 필요에 대한 언급 없이 순전히 형식적 정의로만
설명하는 것이 부적절했는지를 설명해 준다. 아마도 이러한 관점이
법리학(jurisprudence)에 주는 가장 큰 공헌은, 법의 성격에 대한 논의를
자주 흐리게 해왔던 몇몇 오해를 불러일으키는 이분법(dichotomies)들로부터
벗어날 수 있는 여지를 제공한다는 점이다. 예컨대, 전통적인 질문인 ``모든
법체계는 제재를 반드시 포함해야 하는가?''라는 문제는, 자연법의 이 단순한
판본이 제공하는 관점에서 새롭고 보다 명료한 방식으로 제기될 수 있다.
우리는 더 이상 다음의 두 가지 부적절한 대안 사이에서 선택을 강요받지
않아도 된다: 하나는, 법(law)이나 법체계(legal system)라는 단어의
`정의'에 비추어 제재가 요구된다고 주장하는 입장이고, 다른 하나는,
대부분의 법체계가 실제로 제재를 포함하고 있다는 것이 단지 `사실'(just a
fact)에 불과하다고 주장하는 입장이다. 이 두 입장은 모두 만족스럽지 않다.
중앙집중적 제재가 없는 체계에 대해 `법'이라는 용어를 사용하는 것을
금지하는 확립된 원칙은 존재하지 않으며, 제재를 포함하지 않는 체계를
`국제법'(international law)이라 부르는 데에도 타당한 이유(비록 강제성은
없지만)가 있다. 그러나 반대로, 인간 존재의 조건에 부합하는 최소한의
목적들을 수행하려는 국내 체계(municipal system)의 경우에는, 그 안에서
제재가 반드시 차지해야 할 위치를 구분해야 한다. 자연적 사실들과
목적들이, 국내 체계 안에서 제재를 가능하고 필수적인 것으로 만든다는 점을
고려하면, 우리는 이것이 하나의 \emphksans{자연적 필연성}(natural
necessity)이라고 말할 수 있다; 그리고 이와 같은 표현은, 인격(persons),
재산(property), 약속(promises)에 대한 최소한의 보호 형태가 국내법의 필수
구성 요소로 기능한다는 점을 전달하는 데도 필요하다. 우리는 바로 이러한
방식으로 ``법은 어떤 내용도 가질 수 있다''(law may have any content)는
실증주의 테제에 응답해야 한다. 왜냐하면 법뿐만 아니라 많은 다른 사회적
제도들을 적절히 기술하기 위해서는, 정의(definitions)나 통상적인 사실
진술들에 더하여, 제3의 범주의 진술을 위한 자리를 남겨 두어야 한다는
중요한 진실이 있기 때문이다. \textbf{(p.~200)} 즉, 그 진리값이 인간과
그들이 살아가는 세계가 현재 지니고 있는 두드러진 성격들을 계속
유지하는지 여부에 의존하는(contingent) 그러한 진술들이다.

\subsection{3. 법적 유효성과 도덕적
가치}\label{uxbc95uxc801-uxc720uxd6a8uxc131uxacfc-uxb3c4uxb355uxc801-uxac00uxce58}

상호 자제(mutual forbearance)의 체계가 제공하는 보호와 이익은, 법과
도덕의 기초를 이루면서도, 각기 다른 사회에서 매우 상이한 범위의
사람들에게 확장될 수 있다. 일정한 제한을 기꺼이 수용하는 한 어떤 인간
집단에게도 이러한 기본적 보호를 부정하는 것은, 모든 현대 국가가 최소한
언어적으로는 존중을 표하는 도덕과 정의의 원칙에 반하는 일이라는 점은
사실이다. 오늘날 널리 공언되는 도덕적 관점은, 적어도 이러한 근본적인
영역에서만큼은 인간은 동등하게 대우받을 자격이 있으며, 차별적 대우는
단순히 타인의 이익을 근거로 정당화될 수 없다는 관념을 바탕에 깔고 있다.

그러나 한편으로는, 어떤 사회의 법이나 통용되는 도덕이 그들의 범위 내에
있는 모든 사람에게 이러한 최소한의 보호와 이익을 제공해야 할 필요는
없으며, 실제로 종종 그러지 않았다. 노예제 사회에서 지배 집단은 노예를
단순한 사용의 대상이 아니라 인간이라는 사실 자체를 인식하지 못하게 될 수
있으며, 동시에 서로 간의 권리와 이해관계에는 매우 민감한 도덕의식을
유지할 수도 있다. 『허클베리 핀(Huckleberry Finn)』에서, 증기선의
보일러가 폭발해 다친 사람이 있는지를 묻는 말에 허크는 ``아뇨, 그냥
깜둥이 한 명이 죽었어요(No'm: killed a nigger)''라고 답한다. 이때 앤트
샐리가 ``그래도 다행이네, 가끔은 사람도 다치거든(Well it's lucky because
sometimes people do get hurt)''이라고 덧붙이는 장면은, 실제로 인간
사회에서 오래도록 지배적이었던 하나의 도덕 체계를 응축해서 보여준다.
그러한 도덕 체계가 지배하는 사회에서, 허크가 뼈저리게 경험했듯이, 노예에
대해 지배 집단 구성원들 사이에서 자연스럽게 작동하는 상호 관심을 연장
적용하려는 시도는 중대한 도덕적 위반으로 간주되며, 그에 상응하는 도덕적
죄책(moral guilt)의 결과를 초래한다. 나치 독일과 남아프리카 공화국의
사례는, 그 시간이 우리에게 불편할 정도로 가깝다는 점에서 유사한 교훈을
제공한다.

물론 어떤 사회의 법은 때로는 통용되는 도덕보다 앞서기도 했지만,
일반적으로 법은 도덕을 뒤따르며, 노예 살인조차도 공적 자원의 낭비이거나,
노예를 소유한 주인의 재산에 대한 침해로만 간주되기도 한다. 심지어
노예제가 공식적으로 인정되지 않는 사회에서도, 인종, 피부색,
\textbf{(p.~201)} 또는 신조 등의 이유로 차별이 존재할 수 있으며, 이러한
차별은 모든 인간이 타인으로부터 최소한의 보호를 받을 자격이 있다는 점을
법체계나 사회적 도덕이 인정하지 않게 만든다.

이러한 인간사의 고통스러운 사실들은, 어떤 사회가 존속 가능하려면 반드시
\emphksans{일부} 구성원에게는 상호 자제의 체계를 제공해야 하지만, 그것이
반드시 모든 이에게 제공되어야 할 필요는 없다는 점을 명확히 보여준다.
앞서 제재(sanction)의 필요성과 가능성에 대해 논의하면서 강조한 바와
같이, 어떤 규칙 체계가 강제력(force)에 의해 시행되려면, 일정 수의
사람들이 그것을 자발적으로 수용해야 한다. 이러한 자발적 협력은 곧
\emphksans{권위(authority)}를 형성하며, 이는 법과 정부의 강제력을 정초하는
기반이다. 하지만 이렇게 형성된 강제력은 두 가지 주요한 방식으로 사용될
수 있다. 첫째, 그 강제력은 규칙으로부터 보호를 받으면서도 이기적으로
그것을 위반하는 범법자(malefactors)에 대해서만 행사될 수 있다. 둘째, 그
강제력은 일정한 수의 피지배 집단(subject group)을 억압하고 그들을
영구적인 열등한 지위에 고정시키는 데 사용될 수도 있다. 이때 피지배
집단의 규모는 지배 집단(master group)의 연대력(solidarity),
규율(discipline), 그리고 강제 수단의 가용성에 따라 크거나 작을 수
있으며, 피지배 집단의 무력함이나 조직화 능력 부족은 그 억압을 더욱
가능하게 만든다. 이러한 억압을 받는 자들에게는, 그 체계가 그들의 충성을
요구할 만한 어떠한 것도 제공하지 않으며, 오직 두려움의 대상일 뿐이다.
그들은 법체계의 수혜자(beneficiaries)가 아니라 피해자(victims)이다.

이 책의 앞선 장들에서 우리는 법체계의 존재는 하나의 사회적 현상(social
phenomenon)이라는 점을 강조했으며, 그것이 항상 두 가지 측면을 지닌다는
점에 주의를 기울여야 한다고 말했다. 하나는 규칙의 자발적 수용(voluntary
acceptance of rules)에 내재된 태도와 행위이고, 다른 하나는 단순한
복종(obedience)이나 묵인(acquiescence)에 따른 보다 단순한 형태의 태도와
행위이다.

따라서 법을 갖춘 사회는, 그 규칙들을 단지 불복종할 경우 공무원에게서
어떤 처벌이 닥칠지를 예측하게 해주는 신뢰할 만한 지표로만 여기는 것이
아니라, 행위의 기준으로 수용된 규범(accepted standards of
behaviour)으로, 즉 내적 관점(internal point of view)에서 바라보는
사람들을 포함한다. 그러나 동시에 그러한 사회는,
범죄자(malefactors)이거나 단순히 그 체계의 무력한 희생자인 사람들처럼,
이러한 법적 규범들이 오직 물리적 강제력 또는 그 위협을 통해 부과되어야만
하는 이들도 포함한다. 이들은 그러한 규칙들을 가능한 처벌의 원천으로만
인식한다. 이 두 부류 사이의 균형은 \textbf{(p.~202)} 다양한 요소들에
의해 결정된다. 만약 그 체계가 공정하고, 순응(compliance)을 요구하는 모든
이들의 중대한 이해관계(vital interests)를 진정성 있게 고려한다면, 그것은
대부분의 사람들에게 대부분의 시간 동안 충성을 얻고 유지할 수 있으며,
결과적으로 안정적인 체계가 될 수 있다. 반대로, 법 체계가 지배 집단의
이익에만 복무하는 협소하고 배타적인 것이라면, 그 체계는 점점 더
억압적이며 불안정해질 것이고, 잠재된 격변의 위협 아래 놓이게 될 것이다.
이 두 극단 사이에는 다양한 법에 대한 태도 조합들이 존재하며, 종종 하나의
개인 내부에도 이러한 이중성이 공존한다.

이러한 측면에 대한 성찰은 하나의 냉정한 진실을 드러낸다. 의무의 일차
규칙(primary rules of obligation)만이 사회적 통제의 유일한 수단인 단순한
사회 형태에서, 중앙집중적으로 조직된 입법부, 법원, 공무원, 그리고 제재를
갖춘 법의 세계로 나아가는 단계는, 확실한 이익을 가져오는 동시에 일정한
대가를 수반한다. 그 이익이란 변화에 대한 적응성, 확실성, 그리고
효율성으로서, 그 규모는 막대하다. 그러나 그 대가는, 중앙집중적으로
조직된 권력이 자신을 지지하지 않아도 되는 다수의 사람들을 억압하는 데
사용될 수 있다는 위험이며, 이는 일차 규칙의 단순한 체제에서는 발생할 수
없었던 방식이다. 이러한 위험이 실제로 현실화되었고 또 다시 현실화될 수
있기 때문에, 우리가 자연법(Natural Law)의 최소 내용(minimum content)으로
제시한 바를 넘어서, 법이 도덕에 반드시 순응\emphksans{해야 하는}(\emph{must}
conform) 어떤 추가적인 방식이 존재한다는 주장은 매우 신중한 검토를
필요로 한다. 이러한 언명들 가운데 다수는, 법과 도덕 사이의 연관성이 어떤
의미에서 필연적이라고 주장되는지를 명확히 하지 못하거나, 면밀히 살펴보면
그것이 참이고 중요하기는 하나, 법과 도덕 사이의 필연적 연결로 제시하는
것은 극히 혼란스러운 의미에 불과한 것으로 드러난다. 우리는 이 장을
마무리하면서 이러한 주장들의 여섯 가지 형태를 검토할 것이다.

 (i) \emphksans{권력과 권위} (Power and authority) 법 체계는 단순히 인간이
인간을 지배하는 힘에 근거할 수 없고 또한 그래서는 안 되므로, 반드시
도덕적 의무감(moral obligation)이나 법체계의 도덕적 가치에 대한
확신(conviction)에 근거해야 한다는 주장이 자주 제기된다. 우리는 이
책의 앞선 장들에서, 단순히 위협에 의해 뒷받침되는 명령(orders backed
by threats)과 복종의 습관(habits of obedience)만으로는 법체계의 기초와
법적 유효성(legal validity)이라는 개념을 설명하기에 불충분하다는 점을
강조하였다. 우리는 \textbf{(p.~203)} 제6장에서 논증한 바와 같이,
이러한 개념들을 이해하려면 수용된 승인 규칙(accepted rule of
recognition)의 개념이 반드시 필요하며, 본 장에서 본 것처럼, 강제력의
존재 조건은 적어도 일부 사람들이 자발적으로 그 체계에 협력하고, 그
규칙들을 수용(acceptance)하는 것이다. 이 의미에서, 법의 강제력은
그것의 수용된 권위(accepted authority)를 전제로 한다는 것은 사실이다.
그러나 '단지 권력에 기반한 법'과 '도덕적으로 구속력을 가지는 것으로
받아들여진 법'이라는 이분법(dichotomy)은 완전히 포괄하는(exhaustive)
구분이 아니다. 매우 많은 수의 사람들이, 그들이 도덕적으로 구속력을
가진다고 생각하지 않는 법에 의해 강제될 수 있으며, 자발적으로 법
체계를 수용하는 이들조차도 반드시 스스로를 도덕적으로 구속받는
존재라고 인식해야 하는 것은 아니다. 물론 그들이 그렇게 인식할 때 가장
안정적인 체계가 될 것이다. 사실, 사람들이 그 체계에 대한
충성(allegiance)을 보이는 이유는 다음과 같이 다양할 수 있다. 장기적
이익에 대한 계산, 타인에 대한 사욕없는 관심(disinterested interest),
반성 없는 전통적 태도, 혹은 단지 타인들이 그렇게 하기에 그에 따르려는
바람 등이다. 법체계의 권위를 수용하는 이들이 양심을 살펴본 끝에,
도덕적으로는 이를 받아들여서는 안 된다고 판단하면서도, 다양한 이유로
여전히 수용을 지속할 수도 있는 것이다.

이러한 상식적인 사실들이 종종 모호해지는 이유는, 사람들이 법적
의무(legal obligation)와 도덕적 의무(moral obligation)를 동일한 어휘로
표현하기 때문이다. 법체계의 권위를 수용하는 이들은 법을 내적
관점(internal point of view)에서 바라보며, 법이 요구하는 바를 다음과
같은 규범적 언어로 표현한다. 이 규범적 언어는 법과 도덕 모두에서
공유되는 것이다.: ``나는/너는 \textasciitilde\textasciitilde 해야
한다(I/You ought)'', ``나는/그는 \textasciitilde\textasciitilde 해야만
한다(I/he must)'', ``나는/그들은 의무를 지닌다(I/they have an
obligation).'' 그러나 그들은 이러한 표현을 통해, 법이 요구하는 바를
실천하는 것이 도덕적으로 정당하다는 \emphksans{도덕적} 판단에 전념하고 있는
것은 아니다. 물론 아무런 부가 설명이 없다면, 어떤 사람이 자신 또는
타인의 법적 의무를 이렇게 표현할 경우, 그것을 이행하는 데 반대할 도덕적
또는 기타 이유가 없다고 생각한다는 추정(presumption)이 가능할 것이다.
그러나 이로부터, 어떤 사항이 법적으로 유효한 의무(legally obligatory)로
인정되기 위해서는 반드시 도덕적으로도 의무적(morally obligatory)으로
수용되어야 한다는 결론이 도출되는 것은 아니다. 우리가 말한 이 추정은,
어떤 사람이 법적 의무를 인정하거나 언급하는 것이, 만약 그가 그것을
이행해서는 안 된다고 여기는 확고한 도덕적 혹은 다른 이유를 가지고 있다면
무의미할 것이라는 점에 기초한 것이다.

 (ii) \emphksans{도덕이 법에 미치는 영향} (The influence of morality on law) 모든
현대 국가의 법은 \textbf{(p.~204)} 수많은 지점에서, 통용되는 사회적
도덕(social morality)과 보다 보편적인 도덕적 이상(moral ideals)의
영향을 보여준다. 이러한 영향은 입법을 통해 갑작스럽고 공공연하게 법
안으로 들어오기도 하고, 사법절차를 통해 조용히 점진적으로 스며들기도
한다. 어떤 체계들, 예컨대 미국의 경우에는, 법적 유효성(legal
validity)의 궁극적 기준에 정의(justice)나 실질적 도덕 가치(substantive
moral values)가 명시적으로 포함되어 있다. 반면 영국과 같은 체계에서는,
최고 입법기관의 입법권에 공식적인 제한이 없음에도 불구하고, 그 입법은
정의나 도덕에 대한 신중한 부합(conformity)을 결코 덜 요구하지 않는다.
법이 도덕을 반영하는 추가적인 방식들은 무수히 많으며, 아직 충분히
연구되지 않았다. 예컨대, 어떤 법률 조항은 단지 법적 외피(legal
shell)에 불과하고, 그 조항의 명문 자체가 도덕적 원칙의 도움을 통해
내용을 보완할 것을 요구하기도 한다. 집행 가능한 계약(enforceable
contracts)의 범위는 도덕성과 공정성(fairness)의 개념에 따라 제한될 수
있으며, 민사 및 형사상 불법행위에 대한 책임(liability)은 지배적인
도덕적 책임관(moral responsibility)의 관념에 따라 조정되기도 한다.
어떤 실증주의자(positivist)도 이러한 사실들---즉, 법체계의
안정성(stability)은 이와 같은 도덕과의 대응 유형에 부분적으로
의존한다는 사실---을 부인할 수는 없다. 만약 이것이 법과 도덕 사이의
필연적 연결(necessary connection)을 의미하는 것이라면, 그 존재는
인정되어야 한다.

 (iii) \emphksans{해석} (Interpretation) 법은 구체적 사안에 적용되기 위해서는
해석이 요구되며, 사법 절차(judicial processes)의 본질을 흐리는 신화를
현실적 연구로 제거하면, 법의 개방적 직조(open texture)가 창의적 활동의
광대한 영역을 남긴다는 점이 분명해진다. 이러한 창의적 활동은 일부에서
입법적(legislative)이라 불리기도 한다. 선례(precedent)나
성문법(statute)을 해석함에 있어, 판사들은 맹목적이고 자의적인
선택이나, 사전에 고정된 의미를 갖는 규칙들로부터의 `기계적
연역'(mechanical deduction) 사이에만 갇혀 있는 것이 아니다. 종종
그들의 선택은 해석 대상 규칙의 목적이 합당할 것이라는 가정에 의해
이끌린다. 즉, 해당 규칙들은 부정의를 야기하거나 정착된 도덕
원칙(settled moral principles)을 침해하도록 설계된 것이 아니라는
전제이다. 특히 고도의 헌법적 중요성(high constitutional import)이 있는
사안에서의 사법적 판단(judicial decision)은, 하나의 도덕 원칙의 단순한
적용이 아니라, 상이한 도덕적 가치들 사이의 선택을 수반하는 경우가
많다. 법의 의미가 불확실한 상황에서, 도덕이 항상 명확한 해답을
제공한다고 믿는 것은 어리석은 일이다. 이 지점에서 판사들은 다시금
자의적이지도, 기계적이지도 않은 선택을 할 수 있으며, 바로 여기에서
종종 \textbf{(p.~205)} 판사 고유의 미덕(judicial virtues)이 드러난다.
이러한 미덕들이 법적 판단에 적합한 이유는, 이와 같은 활동을 단순히
`입법'이라 부르기를 주저하게 만드는 이유이기도 하다. 이러한 미덕에는
다음이 포함된다: 대안들을 고찰함에 있어의 공정성과 중립성(impartiality
and neutrality), 영향을 받게 될 모든 이들의 이익에 대한 고려, 그리고
수용 가능한 일반 원칙을 근거 있는 결정의 기반으로 삼고자 하는
노력이다. 물론 항상 다양한 원칙들이 존재 가능하기 때문에, 어떤 결정이
유일하게 옳다는 것을 \emphksans{증명}(demonstrate)할 수는 없지만, 그것은
정보에 기초한 공정한 선택의 추론된 산물(reasoned product)로서 수용될
수 있다. 이 모든 과정에서 우리는 경쟁하는 이해관계들 간의 정의 실현을
위한 `저울질'(weighing)과 `균형 조정'(balancing)의 과정을 보게 된다.

이러한 요소들은, 결정을 수용 가능하게 만드는 데 있어 도덕적(moral)이라
불릴 수 있으며, 그것들의 중요성은 거의 누구도 부인하지 않을 것이다.
대부분의 법체계에서 해석을 지배하는 느슨하고 변화 가능한 전통 혹은 해석
규범(canon of interpretation) 역시 이러한 요소들을 대체로 모호하게나마
포함하고 있다. 그러나 이러한 사실들이 법과 도덕 사이의 \emphksans{필연적}
연결(\emph{necessary} connection)을 입증하는 증거로 제시될 경우, 우리는
다음 사실을 기억해야 한다. 즉, 동일한 원칙들이 지켜지는 경우만큼이나
위반되는 경우에도 거의 동일하게 '존중'되어 왔다는 점이다. 오스틴(Austin)
이후 오늘날에 이르기까지, 이러한 요소들이 판단을 이끌\emphksans{어야
한다(should)}고 주장해 온 이들은, 법관의 법 형성(judicial law-making)이
종종 사회적 가치에 대해 맹목적이거나, '자동적'이거나, 이성적 근거가
부족하다는 점을 지적한 비판자들이었다.


 (iv) \emphksans{법에 대한 비판} (The criticism of law) 때로 법과 도덕 사이의
필연적 연결이 존재한다는 주장은, 단지 ``\emphksans{좋은}(good) 법체계는
정의와 도덕의 요구에 일정한 지점에서 부합(conform)해야 한다''는 주장에
불과하기도 하다. 앞 단락에서 언급한 바와 같은 지점들에서 말이다. 일부
사람들은 이것을 자명한 상식(truism)으로 간주할지도 모르지만, 그것은
동어반복(tautology)은 아니다. 실제로 법에 대한 비판 과정에서는, 어떤
도덕적 기준이 적절한지, 그리고 어떤 지점에서의 부합이 요구되는지를
둘러싸고 의견 차이가 있을 수 있다. 즉, 법이 부합해야 할
도덕(morality)이란 해당 집단 내에서 수용된 통념적 도덕인가? 설령
그것이 미신에 기초하거나, 노예나 피지배 계층에게 이익과 보호를
부여하지 않는 경우일지라도? 아니면, 도덕이란 사실에 대한 이성적 믿음에
기초하고, 모든 인간이 동등한 고려와 존중을 받을 자격이 있다고 인정하는
의미에서 계몽된(enlightened) 기준인가?

\textbf{(p.~206)} 법체계가 그 적용 범위에 속하는 모든 인간을 일정한
기본적 보호와 자유를 향유할 자격이 있는 존재로 취급해야 한다는 주장은,
오늘날 법 비판에서 자명하고도 중요한 이상(ideal)으로 일반적으로
받아들여지고 있다. 실제 실천이 이를 따르지 않더라도, 이 이상에 대한
형식적 동의(lip service)는 통상적으로 제공된다. 심지어 모든 인간이
동등한 고려의 권리를 갖는다는 견해를 수용하지 않는 도덕 이론은,
철학적으로 내적 모순, 독단성, 혹은 비이성성을 내포하고 있음이 논증될
수도 있다. 만약 그렇다면, 이러한 권리를 인정하는 계몽된 도덕(enlightened
morality)은 진정한 도덕(true morality)으로서 특별한
정당성(credentials)을 가지며, 단지 가능한 여러 도덕 중 하나가 아니다.
이러한 주장은 여기서 검토할 수 없지만, 비록 이들이 받아들여지더라도, 그
사실은 다음의 사실을 바꾸거나 흐리게 해서는 안 된다. 즉, 기본
규칙(primary rules)과 이차 규칙(secondary rules)의 구조를 특징으로 하는
국내 법체계(municipal legal systems)는, 이러한 정의의 원칙들을
무시하면서도 오랫동안 지속되어 왔다는 점이다. 부정의한 규칙(iniquitous
rules)을 '법'으로 인정하지 않음으로써 무엇을 얻게 되는지에 대한 논의는,
아래에서 다룬다.

 (v) \emphksans{법다움과 정의의 원칙들} (Principles of legality and justice)
도덕성과 정의에 일정한 지점에서 부합하는(conform) 좋은 법체계와 그렇지
않은 법체계 사이의 구분은, 인간 행위가 일반 규칙에 의해 공공연히
공표되고 사법적으로 적용되는 방식으로 통제되는 한, 정의의 최소한이
필연적으로 실현된다는 이유에서, 오류를 포함한 구분이라고 말할 수도
있다. 실제로 우리는 이미 정의(justice)의 개념을 분석하면서\footnote{p.~160
  above.}, 그 가장 단순한 형태인 `법의 적용에서의 정의'(justice in the
application of the law)가 단지 ``서로 다른 많은 사람들에게 적용되어야
할 것이 동일한 일반 규칙(general rule)''이라는 생각을 편견, 이해관계,
또는 변덕에 휘둘리지 않고 진지하게 받아들이는 것에 지나지 않음을
지적한 바 있다. 이러한 공정성(impartiality)은, 영어 및 미국 법률가들이
`자연적 정의'(Natural Justice)의 원칙들로 알고 있는 절차적
기준(procedural standards)이 보장하고자 하는 바로 그 요소이다. 따라서
가장 혐오스러운 법률조차도 정당하게 적용될 수 있으며, 일반 법규칙을
적용한다는 그 자체의 최소한의 관념 속에는 적어도 정의의 씨앗(germ of
justice)이 존재한다고 할 수 있다.

이러한 정의의 최소 형태, 곧 `자연적'(natural)이라 불릴 수 있는 양상들은,
실제로 \textbf{(p.~207)} 어떤 형태의 사회 통제가 작동하려면 어떤
조건들이 전제되어야 하는지를 살펴볼 때 드러난다. 이는 게임의 규칙이든
법이든, 기본적으로는 일반적인 행위 기준(general standards of conduct)을
사람들의 특정 집단에 전달하고, 그 집단이 추가적인 공식 지시 없이 그것을
이해하고 준수(conformity)할 것을 기대하는 방식의 통제이다. 이러한 유형의
사회 통제가 기능하기 위해서는, 규칙은 다음의 조건을 충족해야 한다: 이해
가능해야 하며(intelligible), 대부분의 사람들이 준수할 수 있어야 하고,
일반적으로는 소급 적용(retrospective)이 되어서는 안 된다. 예외적으로만
가능하다. 이 말은, 규칙 위반으로 처벌받는 사람들은 대부분 그것을 준수할
능력과 기회를 가지고 있었다는 것을 의미한다. 명백히 이러한 규칙에 의한
통제의 특징들은, 법률가들이 법다움의 원칙(principles of legality)이라
부르는 정의의 요구사항들과 밀접히 관련되어 있다. 실제로 어떤 실증주의
비판자는, 이러한 규칙에 의한 통제의 양상들 속에서 법과 도덕 사이의
필연적 연결(necessary connection)의 의미를 찾았고, 그것을 `법의 내적
도덕성'(the inner morality of law)이라 불러야 한다고 제안하기도 했다.
다시 말해, 만약 이것이 법과 도덕의 필연적 연결이라는 것이 의미하는
바라면, 우리는 그것을 수용할 수도 있다. 그러나 유감스럽게도, 그것은
극심한 부정의(inquity)와도 양립할 수 있다.

 (vi) \emphksans{법적 유효성(legal validity)과 법에 대한 저항} 그들의 일반적
관점을 다소 조심성 없이 제시했다고 하더라도, 실증주의자로 분류되는
법이론가들 중 앞의 다섯 항목에서 논의한 법과 도덕 사이의 연결 양태들을
부정하고자 했던 사람은 드물다. 그렇다면, 법실증주의의 유명한 구호들이
지향했던 바는 무엇인가? 예컨대, ``법의 존재는 한 가지 문제고, 그
장점(merit) 여부는 또 다른 문제다.''\footnote{Austin, \emph{The
  Province of Jurisprudence Defined}, Lecture V, pp.~184--5.};
``국가의 법은 이상이 아니라 실제로 존재하는 것이다 \ldots{} 그것은
되어야 할 것이 아니라, 있는 그대로의 것이다.''\footnote{Gray,
  \emph{The Nature and Sources of the Law}, s. 213.}; ``법규범(legal
norms)은 어떠한 내용(content)도 가질 수 있다.''\footnote{Kelsen,
  \emph{General Theory of Law and State}, p.~113.}

이러한 사상가들이 주로 옹호하고자 했던 것은, 도덕적으로 부정의하더라도
적법한 형식으로 제정되었고, 의미가 명확하며, 해당 법체계의 모든 인정된
유효성(validity) 기준을 충족하는 특정 법률의 존재가 제기하는 이론적 및
도덕적 문제들을 명료하고 정직하게 서술하려는 시도였다. 그들의 입장은,
이러한 법률에 대해 고찰할 때, 이론가든 그 법률을 적용하거나 준수해야
하는 공무원 또는 시민이든 간에, \textbf{(p.~208)} 그것을 `법' 또는
'유효한 법'이라는 명칭으로 부르지 말라는 초대는 혼란만을 야기할 것이라는
것이었다. 그들은, 이 문제에 직면할 때, 더 단순하고 더 솔직한 방법들이
있으며, 그것이 관련된 모든 지적·도덕적 고려사항들을 더 명확히 조명해 줄
것이라고 생각했다. 즉 우리는 ``이것은 법이다. 그러나 그것은 너무나
부정의하여 적용하거나 복종할 수 없다''라고 말해야 한다는 것이다.

이에 반대되는 입장은, 혁명이나 대규모 정치적 격변 이후에, 법원이 이전
체제하에서 공무원이나 민간인이 법 형식을 통해 저지른 도덕적 불의에 대해
어떤 태도를 취할지를 고찰해야 할 때, 매력적으로 보이기도 한다. 이들의
처벌은 사회적으로 바람직하게 느껴질 수 있으나, 그것을 달성하기 위해서는,
과거 체제 하에서 허용되거나 심지어 요구되었던 행위를 소급적으로
범죄화하는 입법이 필요하며, 이는 실현이 어렵거나, 스스로 도덕적으로
혐오스럽거나, 실행 불가능할 수도 있다. 이러한 상황에서, 법적 어휘 속에
잠재된 도덕적 함의들을 활용하는 것이 자연스럽게 보일 수도 있다. 특히
\emph{ius, recht, diritto, droit}와 같은 단어들은 자연법(Natural Law)의
이론으로 충만한 용어들이므로 더욱 그러하다. 이 경우, 부정의를 명령하거나
허용했던 법률들은, 비록 그것들이 제정된 법체계가 입법권에 대한 어떠한
제한도 인정하지 않았더라도, 유효(valid)하거나 법의 성격(quality of
law)을 갖는 것으로 인정되어서는 안 된다(recognized)는 주장을 하게 되는
유혹에 빠지기 쉽다. 이러한 형태로, 제이차 세계대전 이후 독일에서는 나치
정권의 부정의와 그 패배로 인해 남겨진 첨예한 사회 문제들에 대응하기 위해
자연법 논증들이 부활했다. 이기적 목적을 위해, 나치 정권 하에서 제정된
괴물같은 법령을 위반한 범죄를 저지른 사람들을 투옥시키기 위해 밀고한
사람들은, 처벌받아야 하는가? 그러한 법령들이 자연법을 위반했기 때문에
무효(void)이며, 그로 인해 그 법령을 위반한 것에 따른 투옥은 실질적으로
불법(unlawful)이었고, 그것을 야기한 행위는 그 자체로 범죄였다고 하여,
전후 독일 법정에서 그들을 유죄로 판단할 수 있었는가?\footnote{See the
  judgment of 27 July 1949, Oberlandesgericht Bamberg, 5
  \emph{Süddeutsche Juristen-Zeitung}, 207: discussed at length in H. L.
  A. Hart, `Legal Positivism and the Separation of Law and Morals', in
  71. \emph{Harvard L. Rev.} (1958), 598, and in L. Fuller, `Positivism
  and Fidelity to Law', ibid., p.~630. But note corrected account of
  this judgment below, pp.~303--4.} 이처럼 \textbf{(p.~209)}
``도덕적으로 부정의한 규칙은 법이 될 수 없다''는 견해를 받아들이는
입장과 거부하는 입장 사이의 논쟁은 표면상 간단해 보이지만, 그 일반적
성격에 대한 이해는 종종 불분명하다. 여기서 문제되는 것은, 도덕적으로
부정의한 규칙을 적용하거나, 복종하거나, 타인이 그것을 항변하는 것을
허용하지 않겠다는 도덕적 결정을 표현하는 \emphksans{대안적 방식들} 간의
차이이다. 그러나 이 문제는 단순히 언어적 표현에 관한 문제로 제시될 수는
없다. 논쟁의 양측 누구도, 단순히 ``당신 말이 맞습니다. 당신이 주장한
바는 영어(또는 독일어)에서 그 종류의 주장을 표현하는 올바른
방식입니다''라는 답변에 만족하지 않을 것이다. 따라서 실증주의자는 다음과
같이 주장할 수도 있다. 즉, 영어 사용 관행에 비추어 보아, ``어떤 법규칙이
너무나 부정의하여 복종할 수 없다''고 주장하는 데에는 모순이 없으며, ``그
규칙에 복종할 수 없다''는 명제가 ``그 규칙은 유효한 법규칙이 아니다''는
결론으로 이어지는 것은 아니라고 주장할 수 있다. 그러나 그 반대편에 선
이들은 이 주장으로 인해 논쟁이 종결되었다고 보지는 않을 것이다.

명백히, 우리는 이 문제를 언어 사용의 적절성(propriety)에 관한 것으로만
본다면, 그것을 충분히 다룰 수 없다. 진정으로 문제가 되는 것은,
사회생활에서 일반적으로 효과적으로 작동하는 하나의 규칙 체계에 속한
규칙들을 분류하는 방식으로서의 넓은 개념과 좁은 개념, 즉 경쟁하는 두
규범 개념들 중 어느 것이 이론적 탐구나 도덕적 숙고---혹은 그 둘
모두---를 더 잘 보조하는지에 따라 선택되어야 한다는 점이다.

이 두 경쟁 개념 중에서 더 넓은 개념은 좁은 개념을 포함한다. 더 넓은
개념을 채택할 경우, 우리는 이론적 탐구에서 다음과 같은 방식을 취하게
된다: 일차적(primary) 및 이차적(secondary) 규칙의 체계가 제시하는 형식적
기준에 따라 유효(valid) 한 모든 규칙들을---비록 그 일부가 특정 사회의
도덕성이나 우리가 진정한 도덕성 또는 계몽된 도덕성이라 여기는 것에
어긋나더라도---`법'(law)으로 묶어 함께 고찰하게 된다. 반면, 좁은 개념을
채택하면 그러한 도덕적으로 불의한 규칙들은 '법'의 범주에서 제외된다.
그러나 이처럼 좁은 개념을 채택하는 것은, 법을 사회적 현상으로 연구하는
이론적·과학적 탐구에 있어서 아무런 이익도 가져다주지 않는 듯하다.
왜냐하면 이는 비록 법의 복잡한 특징들을 모두 갖추었더라도, 특정 규칙들을
법에서 배제하게 되기 때문이다. 이러한 규칙들을 다른 학문 분야로 넘기자는
제안에서 비롯될 수 있는 것은 오직 혼란뿐이며, 실제로 그 어떤 역사적
연구나 법학의 다른 분과도 이러한 시도를 유익하다고 판단하지 않았다. 반면
넓은 개념의 법 개념을 채택하면, \textbf{(p.~210)} 우리는 도덕적으로
불의한 법이 지니는 특수한 성격들과 사회가 그것에 보이는 반응까지
포함하여 포괄적으로 연구할 수 있다. 따라서 좁은 개념을 사용할 경우,
일차·이차 규칙 체계에서 구현되는 사회통제의 특정 방식이 어떻게 발전하고
어떤 가능성을 지니는지를 이해하려는 우리의 노력은 혼란스럽게 분열될
수밖에 없다. 그것의 사용에 대한 연구는 필연적으로 그것의 오용(misuse)에
대한 연구를 포함해야 하기 때문이다.

그렇다면 좁은 개념이 도덕적 숙고에 있어 가지는 실천적 유익성은 무엇인가?
도덕적으로 부정의한 요구에 직면했을 때, '이것은 어떤 의미에서도 법이
아니다'라고 판단하는 것이 '이것은 법이지만 너무 부정의하여 복종하거나
적용할 수 없다'라고 판단하는 것보다 나은가? 이 판단은 사람들을 더
명료하게 만들거나, 도덕이 요구할 때 더 쉽게 불복종하도록 만들 것인가?
혹은 나치 정권이 남긴 문제들을 더 나은 방식으로 처리하게 할 것인가? 물론
사상은 영향력을 가진다. 그러나 도덕적으로 부정의하지만 유효(valid) 한
법이 존재할 수 없다는 좁은 법적 유효성(validity) 개념을 사람들에게
교육하고 훈련시키는 노력이, 조직화된 권력의 위협에 직면하여 악에
저항하는 의지를 강화시키거나, 복종이 요구될 때 도덕적으로 어떤 문제가
제기되는지를 더 명확하게 인식하게 만들 것이라 보기는 어렵다. 인간이
일부의 협조를 통해 타인을 지배할 수 있는 한, 그들은 법의 형식(form)을 그
수단 중 하나로 사용할 것이다. 사악한 자들은 사악한 규칙을 제정하고, 다른
자들은 그것을 집행할 것이다. 이러한 권력의 공식적 남용에 맞서 사람들을
명확하게 인식하게 만들기 위해 무엇보다 필요한 것은, 어떤 규칙이 법적으로
유효하다고 인증(certification)되었다는 사실이 복종의 문제를 최종적으로
해결하는 것은 아니라는 점을 인식하는 것이다. 그리고 공식적 체계가 아무리
위엄이나 권위를 띠더라도, 그 요구는 결국 도덕적 심판(moral scrutiny)에
복종되어야 한다는 점을 자각하는 것이다. 이러한 감각---곧 공식적 체계
밖에 어떤 기준이 존재하며, 개인은 최종적으로 그것에 의거해 복종의 문제를
해결해야 한다는 인식---은 법규칙이 부정의할 수 있다(may be iniquitous)는
사고에 익숙한 사람들 사이에서 훨씬 더 생생히 유지될 가능성이 크다. 반면,
어떤 부정의한 것도 결코 법의 지위를 가질 수 없다고 생각하는 사람들
사이에서는 그렇지 않다.

그러나 우리가 `이것은 법이지만 부정의하다'라고 사고하고 표현할 수 있게
해주는 법의 넓은 개념을 선호해야 할 더 강력한 이유는 따로 있다.
\textbf{(p.~211)} 그것은 부정의한 규칙에 대해 법적 인정을 거부함으로써,
그것이 초래하는 도덕적 문제들의 다양성과 복잡성을 지나치게 단순화할
위험이 있기 때문이다. 벤담(Bentham)과 오스틴(Austin)과 같은 고전
이론가들은 `법이란 무엇인가'와 `법은 어떻게 되어야 하는가' 사이의 구분을
강조했는데, 이는 사람들이 이 구분을 유지하지 않으면, 사회에 어떤 대가가
따르든 무시한 채 법이 유효하지 않다고 성급히 판단하고 복종하지 않으려 할
수 있기 때문이었다. 이러한 무정부 상태의 위험은 그들이 과대평가했을 수도
있으나, 또 다른 형태의 과도한 단순화의 위험도 있다. 만약 우리가 관점을
좁혀 '악한 규칙에 \emphksans{복종(obey)}해야 하는 사람'만을 고려한다면, 그가
그것이 '법'이라는 유효한 규칙인지 아닌지를 어떻게 인식하든, 도덕적
부정의를 인식하고 도덕이 요구하는 대로 행동한다면 그것은 무관하다고
판단할 수도 있다. 그러나 '복종'이라는 도덕적 문제(``내가 이 부정의한
일을 해야 하는가?'') 외에도, 소크라테스가 던진 문제---``불복종에 대한
처벌을 감수해야 하는가, 아니면 탈출해야 하는가?''---가 있으며, 또한 전후
독일 법원이 마주했던 질문---``당시 시행 중이던 부정의한 규칙들이 허용한
악행을 저지른 자들을 우리는 처벌해야 하는가?''---도 존재한다. 이러한
질문들은 서로 매우 다른 도덕적·정의적 문제들을 제기하며, 각각 독립적으로
고려되어야 한다. 그 어떤 목적을 위해서도 부정의한 법률을 유효하다고
인정하지 않겠다(refuse to recognize)는 일회적 선언으로 해결될 수 없다.
이것은 복잡하고 미묘한 도덕적 문제들에 대해 지나치게 거친 접근 방식이다.

법의 유효성(validity)을 그 도덕성(morality)으로부터 구별할 수 있도록
해주는 법 개념은, 이처럼 서로 다른 문제들의 복잡성과 다양성을 인식할 수
있게 한다. 반면, 부정의한 규칙은 결코 법적으로 유효할 수 없다고 주장하는
좁은 법 개념은, 오히려 이러한 문제들을 가리는 역할을 할 수도 있다.
독일의 밀고자들이 이기적인 동기에서 괴물 같은 법 아래에서 타인을
처벌받게 한 것이 도덕적으로 잘못된 일임은 분명하지만, 동시에 도덕은
국가는 오직, 그 당시 법이 금한 악행을 저지른 사람만을 처벌해야 한다고
요구할 수도 있다. 이것이 바로 \emph{`법 없이는 형벌도 없다'(nulla poena
sine lege)라는} 원칙이다. 만약 이 원칙을 희생하지 않고서는 더 큰 악을
피할 수 없다면, 그 희생이 정당화되는 조건들이 명확히 규정되어야 한다.
소급 처벌의 사례가, 마치 당시 이미 불법이었던 행위에 대한 처벌인 것처럼
보이게 해서는 안 된다. \textbf{(p.~212)} 실증주의적 간명한 입장이, 즉
'도덕적으로 부정의한 규칙들도 여전히 법일 수 있다'는 주장이 가지는
최소한의 의의는 바로 이것이다. 그것은 극단적 상황에서 악 중의
선택(choice between evils)이 불가피할 때, 그 선택을 덮거나 위장하지
않는다는 점이다.

\subsection{CHAPTER IX 주석}\label{chapter-ix-uxc8fcuxc11d}

\emph{185쪽. 자연법(Natural Law).} 고전적·스콜라철학적·근대적 자연법
개념과 '실증주의(positivism)'라는 표현의 다의성에 대한 방대한 주석
문헌이 존재하는 탓에, 자연법(Natural Law)이 법실증주의(Legal
Positivism)와 대립할 때 정확히 어떤 쟁점이 걸려 있는지를 파악하기 어려울
때가 많다. 본문에서는 그 중 하나의 쟁점을 식별하려는 시도를 한다. 그러나
이 주제를 논의함에 있어 \emphksans{이차 문헌}만을 읽는다면 얻을 수 있는 것은
극히 제한적이다. \emphksans{자연법의 일차 문헌들}에 나타난 어휘와 철학적
전제들을 직접적으로 이해하는 것이 필수적이다. 다음은 쉽게 접근 가능한
최소한의 자료들이다: 아리스토텔레스(Aristotle), 『자연학(Physics)』
제2권 제8장 (Ross 역, Oxford); 아퀴나스(Aquinas), 『신학대전(Summa
Theologica)』 제1-2부, 문제 90--97 (D'Entrèves, 『아퀴나스: 선별된
정치적 저작』(Aquinas: Selected Political Writings), Oxford, 1948에 번역
수록); 그로티우스(Grotius), 『전쟁과 평화의 법(On the Law of War and
Peace)』, 서론(Prolegomena) (The Classics of International Law, 제3권,
Oxford, 1925 수록 번역본); 블랙스톤(Blackstone), 『주석(Commentaries)』,
서론 제2절.

\emph{185쪽. 법실증주의(Legal Positivism).} 현대 영미권 문헌에서
`실증주의(positivism)'라는 표현은 다음과 같은 하나 혹은 그 이상의 주장을
가리키는 데 사용된다: (1) 법은 인간에 의해 부과된 명령(command)이다; (2)
법과 도덕, 또는 존재하는 법(law as it is)과 마땅히 그래야 할 법(law as
it ought to be) 사이에는 필연적 연결이 없다; (3) 법 개념의 의미를
분석하거나 연구하는 것은 역사적 탐구, 사회학적 탐구, 도덕·사회적
목적·기능에 따른 비판적 평가와는 구별되는 (그러나 적대적인 것이 아닌)
중요한 연구 분야이다; (4) 법체계는 `폐쇄된 논리 체계(closed logical
system)'이며, 올바른 판결은 오직 미리 정해진 법규로부터 논리적
방법만으로 도출될 수 있다; (5) 도덕판단은 사실 명제(statements of
fact)처럼 합리적 논변, 증거, 또는 입증을 통해 정립될 수 없다(윤리적
비인지주의, non-cognitivism in ethics). 벤담(Bentham)과 오스틴(Austin)은
(1), (2), (3)의 견해를 지지했으나 (4), (5)는 지지하지 않았다.
켈젠(Kelsen)은 (2), (3), (5)를 수용했으나 (1), (4)는 수용하지 않았다.
주장 (4)는 흔히 `분석법학자(analytical jurists)'에게 귀속되곤 하지만,
이는 타당한 근거 없이 이루어진 것이다. 대륙 유럽 문헌에서는
`실증주의(positivism)'라는 표현이 흔히 \emphksans{이성에 의해서만 발견 가능한
인간 행위의 어떤 원칙이나 규칙의 존재}를 전면적으로 부정하는 입장을
가리키는 데 사용된다. `실증주의' 개념의 다의성에 대한 아고(Ago)의 유익한
논의는 \emph{American Journal of International Law} 제51권(1957) 참조.

\emph{186쪽. 밀(Mill)의 자연법 비판.} 『자연(Nature), 종교의 효용성,
유신론』에 수록된 ``자연(Nature)'' 에세이 참조.

\emph{187쪽. 블랙스톤과 벤담의 자연법에 대한 견해.} 블랙스톤, 위
인용(loc. cit.); 벤담, 『주석에 대한 논평(Comment on the Commentaries)』
제1--6절.

\emph{193쪽. 자연법의 최소한의 내용.} 자연법의 이 경험주의적 버전은
홉스(Hobbes), 『리바이어던(Leviathan)』 제14·15장과 흄(Hume),
『인성론(Treatise of Human Nature)』 제3권 제2부, 특히 제2절 및
제4--7절에 근거한다.

\emph{200쪽. 『허클베리 핀(Huckleberry Finn)』.} 마크 트웨인(Mark
Twain)의 이 소설은 사회적 도덕이 개인의 공감(sympathy)이나
인도주의(humanitarianism)에 반할 경우 발생하는 도덕적 딜레마에 대한
심오한 탐구다. 모든 도덕성을 인도주의와 동일시하는 경향에 대한 귀중한
반론이 된다.

\emph{200쪽. 노예제(Slavery).} 아리스토텔레스에게 있어 노예는 '살아있는
도구'였다. (\emphksans{정치학(Politics)} 제1권 제2--4장).

\emph{203쪽. 도덕성이 법에 미치는 영향.} 법의 \emphksans{발전(development)}이
도덕성에 의해 어떻게 영향을 받아왔는지를 다룬 유익한 연구들로는
에임스(Ames), ``법과 도덕(Law and Morals)'', \emph{Harvard Law Review}
제22권(1908); 파운드(Pound), 『법과 도덕(Law and Morals)』(1926);
굿하트(Goodhart), 『영국법과 도덕법(English Law and the Moral
Law)』(1953) 등이 있다. 오스틴은 이와 같은 사실적 혹은 인과적 연관을
충분히 인식하고 있었다. 『법의 영역(The Province of Jurisprudence
Determined)』 제5강, 162쪽 참조.

\emph{204쪽. 해석(Interpretation).} 법 해석에서 도덕적 고려가 차지하는
위치에 대해서는 라몽(Lamont), 『가치판단(The Value Judgment)』
296--301쪽; 웩슬러(Wechsler), ``헌법 해석의 중립 원칙을 향하여'',
\emph{Harvard Law Review} 제73권, 960쪽; 하트(Hart), 앞서 인용한
\emph{HLR} 제71권, 606--615쪽 및 풀러(Fuller)의 비판, 동일호 661쪽
\emphksans{끝 부분(ad fin.)} 참조. 오스틴의 `경쟁하는 유추(competing
analogies)' 사이에서 판사가 선택할 수 있는 여지를 인정한 부분과,
공리성의 기준에 그 판결이 부합하지 않는다는 비판은 『강의(The
Lectures)』 제37·38강 참조.

\emph{205쪽. 법에 대한 비판과 모든 인간의 평등한 고려에 대한 권리.}
이러한 권리가 단지 다양한 도덕성 중 하나에 속하는 것이 아니라
\emphksans{진정한 도덕성의 정의적 특징(defining feature)}이라는 견해는
벤(Benn)과 피터스(Peters), 『사회 원칙과 민주국가(Social Principles and
the Democratic State)』 제2·5장, 그리고 바이어(Baier), 『도덕적 관점(The
Moral Point of View)』 제8장을 참조.

\emph{206쪽. 법다움과 정의의 원칙(Principles of legality and justice).}
법다움 원칙에 대해서는 홀(Hall), 『형법의 원칙(Principles of Criminal
Law)』 제1장을, '법의 내재적 도덕성(the internal morality of law)'에
대해서는 풀러(Fuller), \emph{Harvard Law Review} 제71권(1958),
644--648쪽 참조.

\emph{208쪽. 전후 독일에서의 자연법 이론의 부활.} G.
라트브루흐(Radbruch)의 후기 견해, 하트(Hart), 그리고 이에 대한
풀러(Fuller)의 응답에 대한 논의는 모두 \emph{HLR} 제71권(1958) 참조. 이
논의는 1949년 7월 바이에른주 고등법원(Oberlandsgericht Bamberg)의 판결을
중심으로 진행되었는데, 이 사건은 1934년 나치 법률 위반 혐의로 남편을
신고한 아내가 남편의 자유를 불법적으로(unlawfully) 박탈한 혐의로 유죄
판결을 받은 사건이다. 이 사건에 대한 설명은 \emph{HLR} 제64권(1951)
1005쪽에 기술된 것으로, 해당 법원이 1934년 법률을 무효로 판단했다는
전제에 따라 이루어진다. 그러나 이 설명의 정확성은 파페(Pappe)의
비판(``나치 시대의 사법판단의 유효성에 대하여'', \emph{Modern Law
Review} 제23권(1960))에 의해 최근에 도전받았다. 파페 박사의 비판은
정당하며, 하트가 논한 사례는 엄밀히 말하면 가상의 것으로 간주되어야
한다. 파페 박사에 따르면(위 저서, 263쪽), 실제 사건에서 항소법원은
법률이 자연법(Natural Law)을 위반한 경우 불법이 될(be unlawful) 수
있다는 이론적 가능성을 받아들이면서도, 쟁점이 된 나치 법률이 이를
위반했다고 보지는 않았다. 피고는 정보를 제공할 \emphksans{의무(duty)}가
없음에도, 순전히 사적 동기에서 그렇게 했으며, 그러한 행위가 '모든 양심
있는 인간의 건전한 양식과 정의감에 반하는 것임을 인식했음이 분명하다'는
이유로 \emphksans{자유의 불법적 박탈(unlawful deprivation of liberty)} 혐의로
유죄 판결을 받았다. 파페 박사가 유사 사건에 대한 독일 연방대법원의
판결을 세밀하게 분석한 내용(268쪽 \emphksans{끝 부분})도 반드시 참고할 필요가
있다.

\subsection{CHAPTER IX 제3판
주석}\label{chapter-ix-uxc81c3uxd310-uxc8fcuxc11d}

\emph{185--6쪽. 실증주의(Positivism)와 `분리가능성 테제(separability
thesis)'.} 하트는 ``법이 도덕의 특정 요구를 재현하거나 충족시키는 것이
필연적 진리인 것은 결코 아니다''라고 말한다. (홈즈 강연에서는 ``법과
도덕 사이에는 필연적 연결이 없다''고 말한다: H. L. A. Hart, ``Positivism
and the Separation of Law and Morals'' (1957) \emph{Harvard Law Review}
제71권 593쪽, 601쪽 주석 25 참조.) 이 \emphksans{분리가능성 테제(separability
thesis)}는 혼란스러운 개념인데, 그 이유 중 하나는 하트 자신이 \emphksans{법과
정의 사이의 규칙적 연관성}(159--60쪽 주석 참조)과 \emphksans{최소한의 내용
명제(minimum content thesis)}(193쪽 주석 참조)를 통해 두 가지 '필연적
연결'을 옹호하는 것처럼 보이기 때문이다. 이 난제를 정리하려는 다양한
시도로는 다음이 있다: John Gardner, 「Legal Positivism: 5½ Myths」,
\emph{Law as a Leap of Faith} 제2장; Matthew H. Kramer, 「On The
Separability of Law and Morality」 (2004) \emph{Canadian Journal of Law
and Jurisprudence} 제17권 315쪽; James Morauta, 「Three Separation
Theses」 (2004) \emph{Law and Philosophy} 제23권 111쪽; Leslie Green,
「Positivism and the Inseparability of Law and Morals」 (2008) \emph{New
York University Law Review} 제83권 1035쪽.

\emph{185--93쪽. 고전적 자연법 이론(Classical theories of natural law).}
하트가 묘사한 견해의 후손 격이나 \emphksans{신토미즘적(neothomist)} 성향이
가미된 입장으로는 John Finnis의 『Natural Law and Natural Rights』, Mark
C. Murphy의 『Natural Law in Jurisprudence and Politics』(Cambridge
University Press, 2006)이 있다. Nigel Simmonds의 『Law as a Moral
Idea』(Oxford University Press, 2007)와 비교할 것. 자연법과 실증주의의
관계에 대해서는 John Finnis, 「The Truth in Legal Positivism」,
『Philosophy of Law』 제7장 참조. 자연법적 도덕 이론, 법 해석 이론,
자연법 법철학 이론 간의 (결여된) 관계에 대해서는 Philip Soper, 「Some
Confusions about Natural Law」 (1992) \emph{Michigan Law Review} 제90권
2393쪽.

\emph{193--200쪽. `최소한의 내용' 명제(The `minimum content' thesis).}
인간의 본성과 우리가 살아가는 세계를 고려할 때, 보편적으로 가치 있는
것들에 관한 \emphksans{최소한의 내용(minimum content)}을 다루지 않는 규칙
체계는 어떤 경우에도 유효한 법체계(legal system)가 될 수 없다. 하트의
목록을 수정한 사례로는 Neil MacCormick, 『H. L. A. Hart』 92--99쪽, John
Finnis, 『Natural Law and Natural Rights』 제4장 참조. 법의 필수적
내용에 대한 다른 접근은 Joseph Raz, 『Practical Reason and Norms』
162--170쪽 참조.

\emph{203쪽. 법과 도덕의 공유 어휘(The shared vocabulary of law and
morality).} `의무(obligation)', `권리(right)', `자유(liberty)' 등의
개념이 법과 도덕에서 공통적으로 사용되는 이유에 대한 다른 설명은 Joseph
Raz, 『The Authority of Law』 제8장 참조. 이에 대한 하트의 응답은
『Essays on Bentham』 153--161쪽에 실려 있다.

\emph{203--4쪽. 도덕의 법 내포(The incorporation of morality in law).}
이 부분, 그리고 더 강조되는 \emphksans{후기(Postscript)}에서 하트는 다음과
같이 말한다: ``일부 법체계에서는 \ldots{} 법적 유효성(legal validity)의
궁극적 기준이 정의 원칙이나 실질적 도덕 가치(substantive moral values)를
명시적으로 내포하고 있다.'' 하트는 이를 `연성 실증주의(soft
positivism)'라 부르지만, 현재는 일반적으로 `내포 명제(incorporation
thesis)' 또는 '포괄적 법실증주의(inclusive legal positivism)'라 불린다.
관련 문헌은 기술적으로 복잡할 수 있다. 훌륭한 개관으로는 K. E. Himma,
「Inclusive Legal Positivism」, Jules Coleman \& Scott Shapiro 편, 『The
Oxford Handbook of Jurisprudence and Philosophy of Law』 수록. 하트의
입장에 대한 비판은 Joseph Raz, 『The Authority of Law』 제3장, 『Between
Authority and Interpretation』 제7장 참조. 포괄적 실증주의를 발전시킨
논의로는 다음이 있다: Jules Coleman, 「Negative and Positive
Positivism」 (1982) \emph{Journal of Legal Studies} 제11권 139쪽; 『The
Practice of Principle: In Defence of a Pragmatist Approach to Legal
Theory』(Oxford University Press, 2001) 제8장; W. J. Waluchow,
『Inclusive Legal Positivism』(Oxford University Press, 1994).
드워킨(Dworkin)은 하트의 접근(Philip Soper와 David Lyons가 채택한
방식)을 『Taking Rights Seriously』 345--350쪽에서 비판하며, Coleman의
입장은 「Thirty Years On」 (2002) \emph{Harvard Law Review} 제115권
1655쪽에서 비판한다.

\emph{206--7쪽. 법다움과 정의의 원칙(Principles of legality and
justice).} 론 풀러(Lon L. Fuller)의 논의가 영향력 있다: 『법의
도덕성(The Morality of Law)』(개정판, Yale University Press, 1969)
46--91쪽; 하트의 서평 「Lon L. Fuller, `The Morality of Law'」는
『Essays in Jurisprudence and Philosophy』 제16장에 수록. 하트와 풀러 간
논쟁의 중요성에 대해서는 Peter Cane 편, 『The Hart--Fuller Debate: 50
Years On』(Hart Publishing, 2010) 수록 논문들 참조. \emphksans{법의 지배(rule
of law)}와 법의 본성 사이의 관계에 대해서는 Raz, 『The Authority of
Law』 제11장 참조. 드워킨은 「Hart's Postscript and the Character of
Political Philosophy」 (2004) \emph{Oxford Journal of Legal Studies}
제24권 1, 23--37쪽에서 '법다움(legality)'이라는 가치에 기초하여 자신의
이론을 재구성한다.

\emph{207--12쪽. 법적 유효성과 법 저항(Legal validity and resistance to
law).} 정의를 확보하려는 최소한의 시도조차 하지 않는 법은 \emphksans{유효하지
않다(invalid)}는 라드브루흐(Gustav Radbruch)의 논제를 현대적으로 계승한
사례로는 Robert Alexy, 『The Argument from Injustice: A Reply to Legal
Positivism』 (Bonnie L. Paulson \& Stanley L. Paulson 번역, Oxford
University Press, 2002). 법이 우리의 복종(obedience)을 요구할 수 있는
도덕적 정당성에 대해서는 다음 문헌 참조: Joseph Raz, 『The Authority of
Law』 제12--15장; A. J. Simmons, 『Moral Principles and Political
Obligations』 (Princeton University Press, 1979); Leslie Green, 『The
Authority of the State』; W. A. Edmundson, 『Three Anarchical Fallacies:
An Essay on Political Authority』 (Cambridge University Press, 1998).

\emph{208--12쪽. '광의의 법 개념(broad concept of law)'을 옹호하기.}
하트는 \emphksans{실증주의적 법 개념}이 ``도덕적 숙고(moral deliberation)를
촉진하고 명확하게 할 수 있다''고 주장한다 (209쪽). 하트는 이 주장을
자신의 분석의 타당성(validity)을 증명하기 위한 필수 요소로 보지는
않으며, 부차적 고려이거나, 광의의 법 개념이 정치적으로 위험하다고
비판하는 사람들에 대한 응답으로 간주한다. 문헌에서는 \emphksans{법 개념(the
concept of law)}이 도덕적으로 바람직한 결과를 위해 조정되어야 한다는
주장과 \emphksans{법의 내용(the content of the law)}이 도덕적으로 좋은 결과를
위해 조정되어야 한다는 주장 사이에 모호함이 존재한다 (예: 판사의 직무를
오직 \emphksans{근원적 자료(source-based materials)} 적용에 한정하려는 시도
등). 비교 자료: Neil MacCormick, 「도덕주의적 법 비판, 비도덕주의적 법
옹호(A Moralistic Case for A-Moralistic Law)」 (1985) \emph{Valparaiso
Law Review} 제20권 1쪽; Thomas Campbell, 『윤리적 실증주의의 법이론(The
Legal Theory of Ethical Positivism)』 (Dartmouth, 1996); Jeremy Waldron,
「규범적(또는 윤리적) 실증주의」, Jules Coleman 편, 『Hart's
Postscript』; Liam Murphy, 「법 개념에 대한 정치적 물음」, 동일 저서
수록.
